\section{1. Overview (Informative)}\label{overview-informative}

This document defines the architectural framework, performance
expectations, and interoperability objectives for \textbf{TimeCard}
devices; modular timing subsystems that provide standardized, high
precision time, phase, and frequency services to a host system.

The primary purpose of this specification is to establish a consistent
structural and logical framework that accommodates various hardware
implementation approaches (e.g., PCIe add-in cards, deeply embedded
silicon IP, or external modules) while facilitating broad industry
interoperability. This interoperability allows diverse TimeCard
implementations to be seamlessly swappable with minimal to no changes
required from the host system, empowering the scalable deployment of
precision timing in hyperscale computing, telecommunications, and
distributed Artificial Intelligence (AI) infrastructure.

As an introductory and informative chapter, the concepts presented
herein provide context for the standard and do not contain normative
conformance requirements.

\begin{center}\rule{0.5\linewidth}{0.5pt}\end{center}

\subsection{1.1 Purpose and Scope}\label{purpose-and-scope}

The \textbf{IEEE P3335 TimeCard Specification} establishes the
structural, electrical, software, and performance characteristics of
TimeCard-based architectures. It describes how hardware subsystems
acquire external time references, maintain temporal stability during
reference loss (holdover), and accurately distribute disciplined phase
and frequency outputs to computing hosts and downstream networks.

The scope of this standardization effort encompasses: - Hardware and
software architecture definitions, including unified timescale
generation and holdover behavior. - Receive (inbound) and Providing
(outbound) timing interface behaviors (e.g., GNSS, PTP, 1PPS, and PCIe
PTM). - Out-of-band and in-band management, telemetry, and security
mechanisms (e.g., I²C, IPMI, REST). - Quantifiable performance metrics
(e.g., ADEV, TDEV, MTIE) and environmental operating constraints. -
Structural methodologies for testing and claiming conformance, including
functional, performance, and environmental validation procedures based
on traceable metrology.

This document is designed to apply equally across: - Vendor-specific
commercial mass-production implementations. - Open-source hardware
reference designs. - Laboratory, exploratory, or research-grade
prototypes.

By defining a \textbf{standardized timing subsystem interface}, the
specification establishes a foundation where different TimeCards from
competing vendors can be easily interchanged, upgraded across
generations, or co-deployed across diverse host platforms while
maintaining stringent functional and timing compatibility.

\begin{center}\rule{0.5\linewidth}{0.5pt}\end{center}

\subsection{1.2 Motivation}\label{motivation}

Modern computing workloads increasingly depend on highly precise and
deterministic hardware timing. Modern applications---such as \textbf{AI
model training clusters, high-frequency financial trading, globally
distributed Spanner-style databases, and 5G/6G cellular
networks}---frequently require sub-microsecond or nanosecond-class
synchronization across thousands of independent nodes.

Historically, achieving this level of timing relied on bespoke,
vendor-specific implementations or tightly coupled architectures that
lacked modularity and scale. The \textbf{TimeCard architecture}
addresses these historical bottlenecks by introducing: - A standardized
hardware abstraction layer and signal footprint (e.g., standardized
1PPS, 10 MHz, and Time of Day bounds).\\
- A unified management and telemetry framework supporting diverse
transports for configuration and observability.\\
- A common timing distribution mechanism capable of bypassing host
operating system latencies using hardware-based timestamping (e.g., PCIe
Precision Time Measurement or hardware PTP).\\
- A clear, falsifiable path for cross-vendor conformance certification,
grounded in traceable metrology and repeatable test procedures.

The ultimate motivation is to foster an open, competitive ecosystem for
precision time distribution.

\begin{center}\rule{0.5\linewidth}{0.5pt}\end{center}

\subsection{1.3 Design Philosophy}\label{design-philosophy}

The P3335 TimeCard specification is founded on several core engineering
principles:

{\def\LTcaptype{none} % do not increment counter
\begin{longtable}[]{@{}
  >{\raggedright\arraybackslash}p{(\linewidth - 2\tabcolsep) * \real{0.4800}}
  >{\raggedright\arraybackslash}p{(\linewidth - 2\tabcolsep) * \real{0.5200}}@{}}
\toprule\noalign{}
\begin{minipage}[b]{\linewidth}\raggedright
Principle
\end{minipage} & \begin{minipage}[b]{\linewidth}\raggedright
Description
\end{minipage} \\
\midrule\noalign{}
\endhead
\bottomrule\noalign{}
\endlastfoot
\textbf{Modularity} & The TimeCard operates as a self-contained
subsystem that fully abstracts the complexities of oscillator
disciplining away from the host CPU. \\
\textbf{Interoperability} & Management and timing interfaces are
standardized to facilitate plug-and-play architectural compatibility
across different server fleets and boundary clocks. \\
\textbf{Scalability} & The architecture gracefully scales from
single-node isolated edge servers up to globally synchronized hyperscale
network domains relying on unified ensemble clocks. \\
\textbf{Determinism} & Hardware-assisted generation and timestamping
provide highly predictable, low-jitter, and low-latency outputs, immune
to software scheduling variations. \\
\textbf{Traceability} & Subsystem logic supports establishing an
unbroken chain of measurements reflecting recognized international
standards (e.g., UTC, TAI), measured using calibrated metrology
equipment. \\
\textbf{Security and Integrity} & Remote management channels and
firmware provisioning protocols rely on modern authenticated and
cryptographically verifiable mechanisms, including provisions for data
sanitization. \\
\end{longtable}
}

\begin{center}\rule{0.5\linewidth}{0.5pt}\end{center}

\subsection{1.4 Relationship to Other
Standards}\label{relationship-to-other-standards}

This specification aligns with and periodically references multiple
established global timing and measurement standards. Key related
standards include:

\begin{itemize}
\tightlist
\item
  \textbf{IEEE Std 1588™-2019 (PTPv2.1):} Precision Time Protocol for
  generalized network synchronization.
\item
  \textbf{IEEE Std 802.1AS™-2020:} Timing and Synchronization for
  Time-Sensitive Applications.
\item
  \textbf{ITU-T G.810 / G.8260 / G.8271:} Definitions and boundary
  parameters for carrier-grade synchronization networks.
\item
  \textbf{IEEE Std 1139™ / IEEE Std 1193™:} Standard definitions and
  environmental testing methodologies for fundamental frequency
  metrology quantities (ADEV, TDEV, MTIE).
\item
  \textbf{PCI-SIG PCIe Base Specification:} Specifically outlining
  Precision Time Measurement (PTM) for in-band hardware timestamping.
\item
  \textbf{Open Compute Project (OCP) Time Appliances Project (TAP):} The
  foundational open-hardware consortium that incubated the original
  TimeCard concept.
\end{itemize}

Where applicable, this document cites these standards normatively to
maintain rigid technical consistency across the industry.

\begin{center}\rule{0.5\linewidth}{0.5pt}\end{center}

\subsection{1.5 Structure of the
Specification}\label{structure-of-the-specification}

The standard is organized sequentially into the following major clauses:

{\def\LTcaptype{none} % do not increment counter
\begin{longtable}[]{@{}
  >{\raggedright\arraybackslash}p{(\linewidth - 2\tabcolsep) * \real{0.4167}}
  >{\raggedright\arraybackslash}p{(\linewidth - 2\tabcolsep) * \real{0.5833}}@{}}
\toprule\noalign{}
\begin{minipage}[b]{\linewidth}\raggedright
Chapter
\end{minipage} & \begin{minipage}[b]{\linewidth}\raggedright
Description
\end{minipage} \\
\midrule\noalign{}
\endhead
\bottomrule\noalign{}
\endlastfoot
\textbf{1. Overview} & Provides the motivation, high-level scope, and
conceptual structure (Informative). \\
\textbf{2. Normative References} & Lists the foundational standards
essential for implementing this specification. \\
\textbf{3. Definitions, Acronyms and Abbreviations} & Clarifies exact
terminology and abbreviations utilized throughout the text. \\
\textbf{4. Conformance} & Outlines the formal terminology and
methodology for claiming implementation compliance. \\
\textbf{5. Architecture} & Defines the core TimeCard blueprint,
oscillator interactions, and boundary definitions. \\
\textbf{6. Performance Specifications} & Details the mathematical
metrics and methodologies used to quantify stability and accuracy. \\
\textbf{7. Timing Interfaces} & Specifies the electrical and logical
conduits for receiving and providing synchronization. \\
\textbf{8. Control Interfaces} & Defines the data structures required
for out-of-band and in-band management. \\
\textbf{9. Environment} & Establishes baselines for physical
survivability, thermal tolerance, and RF emissions. \\
\textbf{10. Applications and Best Practices} & Offers structural
guidance on real-world deployment and operational optimization
(Informative). \\
\textbf{Annex A. Metrics} & Provides definitions and measurement
methodologies for timing and frequency metrics (Informative). \\
\textbf{Annex B. Test Procedures} & Provides standardized procedures for
functional, performance, and environmental testing (Informative). \\
\textbf{Annex C. Bibliography} & Lists informational references and
supplementary reading. \\
\end{longtable}
}

\begin{center}\rule{0.5\linewidth}{0.5pt}\end{center}

\subsection{1.6 Intended Audience}\label{intended-audience}

This specification is drafted distinctly for: - Hardware, software, and
systems engineers developing TimeCard-compatible commercial products or
reference designs.\\
- Datacenter architects and system integrators designing synchronized
infrastructure and planning for scale-out timing topologies.\\
- Network operators and site-reliability engineers managing
latency-sensitive distributed environments who require granular
telemetry.\\
- Scientists and researchers requiring high-fidelity timestamping for
metrology and specialized physics experimentation. - Test and validation
engineers utilizing the defined procedures to empirically grade hardware
performance and claim conformance. - Affiliated standards bodies (e.g.,
ITU-T or OCP) seeking structural harmonization with modern end-system
timing frameworks.

Readers are expected to possess a foundational understanding of
frequency control, phase-locked loops, phase noise characteristics, and
packet-based synchronization (e.g., PTP, NTP).

\begin{center}\rule{0.5\linewidth}{0.5pt}\end{center}

\subsection{1.7 Strategic Goals}\label{strategic-goals}

The widespread adoption of the TimeCard architecture aims to achieve the
following outcomes: - \textbf{Broad vendor-neutral interoperability}
among specialized timing payloads and generic host computers, enabling a
mature, commercial off-the-shelf (COTS) ecosystem.\\
- \textbf{Drastically improved absolute precision} within distributed
databases, cellular networks, and distributed AI compute clusters by
driving the timescale directly to the hardware bus (e.g., PCIe).\\
- \textbf{Lowered integration friction and cost} resulting from strictly
defined hardware abstraction interfaces and deterministic holdover
behaviors.\\
- \textbf{Enhanced operational observability} and traceability to
support stringent financial, telecommunications, and governmental
regulatory frameworks.\\
- \textbf{Consistent metrology evaluations} resulting from standardized,
falsifiable performance testing and conformance guidelines.

\begin{center}\rule{0.5\linewidth}{0.5pt}\end{center}

\subsection{1.8 Document Status and
Origin}\label{document-status-and-origin}

This architectural framework heavily draws its foundations from the
pioneering work accomplished within the \textbf{Open Compute Project
(OCP) Time Appliances Project (TAP)}. The transition of this
architecture into the formal \textbf{IEEE P3335} standardization process
reflects the industry's demand for a rigorously governed, universally
recognized timing hardware standard.

Contributors, silicon vendors, and system integrators are actively
encouraged to submit empirical performance data, interoperability
reports, and technical proposals to inform the ongoing evolution of this
specification.

\section{2. Normative References}\label{normative-references}

The following documents are referred to in the text in such a way that
some or all of their content constitutes requirements of this document.
For dated references, only the edition cited applies. For undated
references, the latest edition of the referenced document (including any
amendments or corrigenda) applies.

\emph{Note: In IEEE drafting style, this section only contains
references that are indispensable for the application of the document
(normative). References that are supplementary, educational, or optional
are located in the Informative Bibliography (see Section 2.2).}

\subsection{2.1 Normative References}\label{normative-references-1}

\begin{itemize}
\tightlist
\item
  \textbf{IEEE Std 1139™}, \emph{IEEE Standard Definitions of Physical
  Quantities for Fundamental Frequency and Time Metrology}.
\item
  \textbf{IEEE Std 1193™}, \emph{IEEE Guide for Measurement of
  Environmental Sensitivities of Standard Frequency Generators}.
\item
  \textbf{IEEE Std 1588™-2019}, \emph{IEEE Standard for a Precision
  Clock Synchronization Protocol for Networked Measurement and Control
  Systems} (Precision Time Protocol / PTPv2.1).
\item
  \textbf{IEEE Std 802.1AS™-2020}, \emph{IEEE Standard for Local and
  Metropolitan Area Networks---Timing and Synchronization for
  Time-Sensitive Applications}.
\item
  \textbf{IETF RFC 5905}, \emph{Network Time Protocol Version 4:
  Protocol and Algorithms Specification}.
\item
  \textbf{ITU-T Recommendation G.810}, \emph{Definitions and terminology
  for synchronization networks}.
\item
  \textbf{ITU-T Recommendation G.8260}, \emph{Definitions and
  terminology for synchronization in packet networks}.
\item
  \textbf{ITU-T Recommendation G.8271}, \emph{Time and phase
  synchronization aspects of telecommunication networks}.
\item
  \textbf{PCI-SIG PCIe Base Specification 5.0} (or later), specifically
  citing the \emph{Precision Time Measurement (PTM)} subclauses.
\item
  \textbf{OCP Time Appliances Project (TAP)}, \emph{TimeCard
  Specification}, Open Compute Project Foundation.
\end{itemize}

\begin{center}\rule{0.5\linewidth}{0.5pt}\end{center}

\subsection{2.2 Bibliography (Informative
References)}\label{bibliography-informative-references}

The bibliography lists documents that are referenced for informative,
educational, or background purposes within the text. Their
implementation is not required to claim conformance to this standard.

\emph{Note: In the final compiled IEEE document, this section will
likely be moved to the final Annex (e.g., Annex A: Bibliography).}

\begin{itemize}
\tightlist
\item
  \textbf{CERN / GSI White Rabbit}, \emph{White Rabbit Specification and
  Documentation} (Sub-nanosecond Ethernet synchronization).
\item
  \textbf{IEC 61000-6-x Series}, \emph{Electromagnetic compatibility
  (EMC) - Generic standards limit for emission and immunity}.
\item
  \textbf{IEEE Communications Surveys \& Tutorials},
  \emph{Multi-constellation GNSS receiver performance studies}.
\item
  \textbf{ISO/IEC 17025}, \emph{General requirements for the competence
  of testing and calibration laboratories}.
\item
  \textbf{ITSF Proceedings}, \emph{International Telecom Sync Forum
  (ITSF) - Industrial case studies and best practices}.
\item
  \textbf{NIST Special Publication 1065}, W.J. Riley, \emph{Handbook of
  Frequency Stability Analysis} (2008).
\item
  \textbf{NIST Technical Note 1952}, \emph{Introduction to PTP and
  synchronization techniques}.
\item
  \textbf{NTP.org}, \emph{NTP Reference Implementation (ntpd)
  open-source software project}.
\item
  \textbf{OCP Github Repository}, \emph{Open Compute Project TimeCard
  Hardware and Firmware Repositories} (Open-source schematics and HDL).
\item
  \textbf{OCP TAP Interoperability Whitepaper}, \emph{Framework for
  multi-vendor validation and testing definitions}.
\item
  \textbf{OCP TAP Wiki}, \emph{Community resources, meeting notes, and
  specification archives}.
\item
  \textbf{Turba Technologies}, \emph{Turba Analytics Documentation - AI
  workload and distributed timing analytics platform}.
\end{itemize}

\begin{center}\rule{0.5\linewidth}{0.5pt}\end{center}

\textbf{End of Chapter -- Normative References}

\section{3. Definitions, Acronyms, and
Abbreviations}\label{definitions-acronyms-and-abbreviations}

This chapter defines the specialized terms, acronyms, and abbreviations
utilized throughout the IEEE P3335 specification.

\subsection{3.1 Definitions}\label{definitions}

For the purposes of this document, the following terms and definitions
apply. The \emph{IEEE Standards Dictionary Online} should be consulted
for terms not defined in this clause.

\begin{itemize}
\tightlist
\item
  \textbf{Accuracy:} The degree of conformance or closeness of a
  measured time or frequency value to a defined primary reference (e.g.,
  UTC).
\item
  \textbf{Asymmetry Calibration:} The process of calculating and
  correcting unequal signal propagation delays in bi-directional timing
  links.
\item
  \textbf{Control Plane:} The logical communication path dedicated to
  configuration, telemetry, and operations management rather than active
  time distribution.
\item
  \textbf{Data Plane:} The logical or physical communication path
  strictly responsible for carrying phase, frequency, and time signals
  to or from the host or network.
\item
  \textbf{Data Plane Latency:} The deterministic or variable propagation
  delay between a time source and its destination hardware.
\item
  \textbf{Deterministic Behavior:} A system characteristic where the
  timing response under specified operational conditions is highly
  predictable and repeatable.
\item
  \textbf{Disciplining:} The continuous process of steering a local
  oscillator's frequency and phase to align with a superior external
  reference signal.
\item
  \textbf{Ensemble Clock:} A highly stable composite time source created
  by mathematically combining multiple independent references via
  weighting or consensus algorithms.
\item
  \textbf{Environmental Chamber:} A specialized testing enclosure used
  to strictly control ambient temperature, humidity, and atmospheric
  conditions for hardware thermal testing.
\item
  \textbf{Frequency Counter:} A measurement instrument used to determine
  the exact signal frequency or period over a defined averaging
  interval.
\item
  \textbf{Granularity:} The minimum distinguishable or settable unit of
  change in a time-domain or frequency-domain output.
\item
  \textbf{Hitless Switching:} The ability of a timing device to execute
  a reference failover or changeover without inducing a measurable phase
  discontinuity to the downstream consumers.
\item
  \textbf{Holdover:} A mode of operation where a timing subsystem
  continues to maintain accurate time based solely on the historical
  stability of its internal oscillator after losing its external
  reference.
\item
  \textbf{Holdover Error:} The progressively increasing deviation or
  drift of the subsystem's timescale from an ideal reference while
  operating in a holdover state.
\item
  \textbf{Host System:} The overarching computing, telecommunications,
  or networking platform that integrates a TimeCard through a
  standardized interface boundary.
\item
  \textbf{Loop Bandwidth:} The effective frequency range over which a
  disciplining Phase-Locked Loop (PLL) tracks the reference; determining
  how fast the loop responds to changes.
\item
  \textbf{Lock Time:} The duration required for the subsystem to acquire
  a reference and assert phase/frequency lock after a cold startup or
  reference transition.
\item
  \textbf{Management Interface:} The designated control channel (e.g.,
  SMBus, IPMI, REST) used strictly for configuration, diagnostics, and
  firmware management.
\item
  \textbf{Oscilloscope (TIE Mode):} Measurement equipment configured to
  capture the Time Interval Error (TIE) of repetitive high-frequency
  signals.
\item
  \textbf{Phase Alignment:} The absolute difference in timing between
  the corresponding edges of two separate signals; typically measured in
  nanoseconds or picoseconds.
\item
  \textbf{Power Analyzer:} An instrument deployed to measure precise
  current draw, power sequencing, and consumption events over time.
\item
  \textbf{Precision:} The statistical repeatability, resolution, or
  variance of a measurement under identical operating conditions.
\item
  \textbf{Primary Reference Clock:} The highest-stability, autonomous
  time or frequency source in a synchronization hierarchy (as defined by
  ITU-T G.811).
\item
  \textbf{Providing Interface:} The outbound physical or logical
  mechanism through which the TimeCard distributes precise time, phase,
  and frequency representations to the host or downstream devices.
\item
  \textbf{Receive Interface:} The inbound physical or logical mechanism
  through which the TimeCard acquires an external reference
  representation.
\item
  \textbf{Reference Signal:} Any external, traceable timing input
  actively utilized to discipline the local oscillator.
\item
  \textbf{Resolution:} The absolute smallest incremental step,
  quantization, or granularity of a time or frequency measurement that
  the system can distinguish.
\item
  \textbf{Rubidium Oscillator:} A highly stable atomic oscillator
  utilizing rubidium-87 vapor transitions to provide exceptional
  long-term frequency stability.
\item
  \textbf{TimeCard:} A modular, hardware-based timing subsystem that
  abstracts synchronization complexity to deliver standardized phase,
  frequency, and time services to a host system.
\item
  \textbf{Time Interval Counter (TIC):} A highly precise laboratory
  instrument designed to measure time differences between distinct
  electrical signal edges with picosecond-level resolution.
\item
  \textbf{Time Jitter:} The short-term variation or instability of a
  time-domain signal (such as a 1PPS edge) from its ideal position,
  frequently reported as an RMS or peak-to-peak value.
\item
  \textbf{Traceability:} A documented, unbroken mathematical chain of
  calibrations linking a local measurement back to recognized primary
  international standards (e.g., UTC).
\item
  \textbf{Unified Timescale:} A single internal temporal scale actively
  maintained by the TimeCard, from which all separate outbound
  interfaces derive coherent, phase-aligned realizations.
\end{itemize}

\begin{center}\rule{0.5\linewidth}{0.5pt}\end{center}

\subsection{3.2 Acronyms and
Abbreviations}\label{acronyms-and-abbreviations}

\begin{itemize}
\tightlist
\item
  \textbf{1PPS:} One Pulse Per Second
\item
  \textbf{ADEV:} Allan Deviation
\item
  \textbf{ASIC:} Application Specific Integrated Circuit
\item
  \textbf{BMC:} Baseboard Management Controller
\item
  \textbf{CXL:} Compute Express Link
\item
  \textbf{CSAC:} Chip-Scale Atomic Clock
\item
  \textbf{DDS:} Direct Digital Synthesizer
\item
  \textbf{EMC:} Electromagnetic Compatibility
\item
  \textbf{EMI:} Electromagnetic Interference
\item
  \textbf{ESD:} Electrostatic Discharge
\item
  \textbf{FPGA:} Field Programmable Gate Array
\item
  \textbf{GNSS:} Global Navigation Satellite System (e.g., GPS, Galileo,
  GLONASS)
\item
  \textbf{gRPC:} Google Remote Procedure Call
\item
  \textbf{I2C:} Inter-Integrated Circuit
\item
  \textbf{I3C:} Improved Inter-Integrated Circuit
\item
  \textbf{IPMI:} Intelligent Platform Management Interface
\item
  \textbf{IRIG:} Inter-Range Instrumentation Group
\item
  \textbf{MAC:} Media Access Control
\item
  \textbf{MMIO:} Memory-Mapped Input/Output
\item
  \textbf{MTBF:} Mean Time Between Failures
\item
  \textbf{MTIE:} Maximum Time Interval Error
\item
  \textbf{NC-SI:} Network Controller Sideband Interface
\item
  \textbf{NTP:} Network Time Protocol
\item
  \textbf{OCP TAP:} Open Compute Project Time Appliances Project
\item
  \textbf{OCXO:} Oven-Controlled Crystal Oscillator
\item
  \textbf{PCIe:} Peripheral Component Interconnect Express
\item
  \textbf{PHY:} Physical Layer
\item
  \textbf{PLL:} Phase-Locked Loop
\item
  \textbf{PN:} Phase Noise
\item
  \textbf{PTM:} Precision Time Measurement
\item
  \textbf{PTP:} Precision Time Protocol
\item
  \textbf{REST:} Representational State Transfer
\item
  \textbf{RoHS:} Restriction of Hazardous Substances
\item
  \textbf{SCPI:} Standard Commands for Programmable Instruments
\item
  \textbf{SEC:} Securities and Exchange Commission
\item
  \textbf{SMA:} SubMiniature version A (connector)
\item
  \textbf{SMBus:} System Management Bus
\item
  \textbf{SNMP:} Simple Network Management Protocol
\item
  \textbf{SWaP-C:} Size, Weight, Power, and Cost
\item
  \textbf{TAI:} International Atomic Time
\item
  \textbf{TCXO:} Temperature-Compensated Crystal Oscillator
\item
  \textbf{TDEV:} Time Deviation
\item
  \textbf{TIC:} Time Interval Counter
\item
  \textbf{TIE:} Time Interval Error
\item
  \textbf{TPM:} Trusted Platform Module
\item
  \textbf{ToD:} Time of Day
\item
  \textbf{USB:} Universal Serial Bus
\item
  \textbf{UTC:} Coordinated Universal Time
\item
  \textbf{WiWi:} Wireless two-Way interferometry
\item
  \textbf{WR:} White Rabbit
\item
  \textbf{WWVB:} Radio Station WWVB (NIST)
\end{itemize}

\begin{center}\rule{0.5\linewidth}{0.5pt}\end{center}

\textbf{End of Chapter -- Definitions, Acronyms, and Abbreviations}

\begin{quote}
\textbf{Editor's Note (Rodney C., P3335 Working Group):} This document
currently serves as a structural guideline for defining what will
eventually go into the Conformance clause of the P3335 standard. The
intention is to align the working group on how to draft normative
requirements.
\end{quote}

\section{4. Conformance Guidelines \&
Methodology}\label{conformance-guidelines-methodology}

\subsection{4.1 IEEE-SA Terminology
Standards}\label{ieee-sa-terminology-standards}

IEEE-SA strictly regulates the terminology used to express requirements
within a standard. The following boilerplate text is provided by the
IEEE-SA template and cannot be edited by the Working Group.

Requirements placed upon conformant implementations of this standard are
expressed using the following specific keywords:

\begin{enumerate}
\def\labelenumi{\arabic{enumi}.}
\tightlist
\item
  \textbf{Shall:} Used exclusively to indicate \textbf{mandatory
  requirements} (must-have features to claim conformance).
\item
  \textbf{May:} Used exclusively to describe \textbf{implementation or
  administrative choices} (``may'' means ``is permitted to'').
  \emph{Note that ``may'' and ``may not'' mean precisely the same thing
  in this context.}
\item
  \textbf{Should:} Used exclusively for \textbf{recommended choices}.
  The behaviors described by ``should'' and ``should not'' are both
  permissible, but one is clearly more desirable than the other.
\end{enumerate}

\subsubsection{4.1.1 Prohibited and Deprecated
Terminology}\label{prohibited-and-deprecated-terminology}

\begin{itemize}
\tightlist
\item
  \textbf{Must:} IEEE-SA strictly prohibits the use of ``must'' to
  describe requirements because it is frequently confused with
  ``shall''. The P3335 text \textbf{cannot} use ``must.'' (If an
  unavoidable physical limitation exists, ``must'' may sometimes be
  justified, but it is generally deprecated in modern drafting).
\item
  \textbf{Will / Can:} These words are used only for statements of fact
  or capability, completely devoid of conformance implications.
\end{itemize}

These three core keywords (\emph{shall, should, may}) are critical. They
instruct an engineer exactly how to implement the standard. Without the
strategic use of these words, the document functions merely as a
``Recommended Practice'' rather than a true Standard.

\begin{center}\rule{0.5\linewidth}{0.5pt}\end{center}

\subsection{4.2 Structural Methodologies for
Conformance}\label{structural-methodologies-for-conformance}

Within IEEE standards, there are three primary methodologies for
organizing and expressing these requirements.

\subsubsection{4.2.1 Methodology A: Conformance Clause
Only}\label{methodology-a-conformance-clause-only}

In this approach, normative clauses (Clause 4 and up) and normative
annexes avoid using the three keywords entirely, instead using
descriptive verbs like ``is'' and ``are'' to explain what an implementer
does. A single, dedicated Conformance Clause at the beginning uses the
keywords by referencing the subsequent normative clauses (e.g.,
\emph{``An implementation of this standard shall implement only one of
the following: a) Clause 4, b) Clause 6, c) Subclause 8.2''}). *
\textbf{Example:} IEEE 802.1

\textbf{Advantages:} * The Conformance Clause serves as a ``one-stop
shop'' to quickly understand the core requirements. * It easily supports
a hierarchical tree of requirements and options (e.g., \emph{``If you
implement high-level option A, you shall implement subclause 4.5''}).

\textbf{Disadvantages:} * Standards authors frequently slip and use
``is'' for concepts unrelated to conformance, which causes ambiguity. *
Authors inevitably accidentally use the three keywords within the
normative text anyway, creating contradictions with the central
Conformance Clause.

\subsubsection{4.2.2 Methodology B: Interspersed Without Conformance
Clause}\label{methodology-b-interspersed-without-conformance-clause}

In this approach, the normative clauses and annexes directly use the
three keywords (\emph{shall, should, may}). Words like ``is'' and
``can'' are reserved solely for factual statements with no conformance
implication. There is no central Conformance Clause; instead, the
requirement tree is reflected purely through the structural organization
of the document (e.g., Subclause 6.8.3 might be explicitly titled
\emph{``Feature foobar (Optional)''}). * \textbf{Example:} IEEE 1588

\textbf{Advantages:} * Avoids the ambiguities and accidental
contradictions of Methodology A. * Conformance is immediately clear as
the implementer reads through the technical text.

\textbf{Disadvantages:} * Places a heavier burden on the implementer to
hunt through the entire document to find all relevant ``shalls.'' * It
is significantly more difficult for authors to document a complex,
mutually exclusive tree of dependencies.

\subsubsection{4.2.3 Methodology C: Interspersed With Conformance
Clause}\label{methodology-c-interspersed-with-conformance-clause}

This is a hybrid approach. It uses the dispersed keywords of Methodology
B but adds a high-level Conformance Clause to summarize the major
requirement trees.

The Conformance Clause provides the high-level roadmap (e.g., \emph{``If
you implement high-level option A, you shall implement subclause
4.5''}). Within subclause 4.5 itself, the text will state things like,
\emph{``The port shall transmit\ldots{}''} which implicitly means
\emph{``If you support 4.5, you shall do this.''}

\begin{center}\rule{0.5\linewidth}{0.5pt}\end{center}

\subsection{4.3 P3335 Working Group
Recommendations}\label{p3335-working-group-recommendations}

\subsubsection{4.3.1 Recommended Drafting
Methodology}\label{recommended-drafting-methodology}

\begin{quote}
\textbf{Recommendation (Rodney C.):} Proceed with drafting the
subsequent text using the \textbf{Interspersed Technique} (Methodology
B).
\end{quote}

As the P3335 project matures, the Working Group can evaluate whether a
high-level Conformance Clause (upgrading to Methodology C) is necessary.
Because P3335 is likely to develop a complex feature tree, a summarizing
Conformance Clause will likely prove highly beneficial in later drafts.

\subsubsection{4.3.2 Protocol Implementation Conformance Statement
(PICS)}\label{protocol-implementation-conformance-statement-pics}

\begin{quote}
\textbf{Recommendation (Rodney C.):} Strongly avoid creating a PICS
annex.
\end{quote}

A PICS typically takes all the dispersed requirements and repeats them
in a massive hierarchical table. Some engineers prefer reading tables
over text for compliance checking, which is why PICS annexes exist in
many IEEE standards.

\textbf{The critical disadvantage:} A PICS \emph{will} inevitably fall
out of sync with the main text of the standard during drafting and
iterative revisions. This creates massive contradictions for
implementers. Adding disclaimers like \emph{``When there is a
contradiction, ignore the PICS''} is often insufficient. \emph{(Note:
IEEE 802.1Q heavily suffers from this exact PICS desynchronization
problem).}

\section{TimeCard Architecture
Specification}\label{timecard-architecture-specification}

The establishment of a standard architecture for TimeCards plays a
critical role in enabling interoperability among diverse
implementations. By defining a consistent framework, different vendors
can design and manufacture TimeCards with varying capabilities,
performance levels, and core technologies, while maintaining full
plug-and-play compatibility with any compliant host. This
standardization fosters an open ecosystem, simplifies hardware
integration, and enables seamless substitution or generational upgrades
of TimeCards without requiring significant host system redesign.

\subsection{5.1 - Architecture Overview}\label{architecture-overview}

A \textbf{TimeCard} is a modular subsystem designed to interface with a
computing host system through a standardized hardware and software
interface. Its primary purpose is to deliver a stable, accurate,
deterministic, and reliable source of time (in the form of time of day,
phase, frequency, or any combination thereof) to the host system.

\subsubsection{5.1.1 Rationale for Dedicated Timing
Subsystems}\label{rationale-for-dedicated-timing-subsystems}

Modern host systems (such as high-performance servers, edge compute
nodes, and telecommunications routers) typically lack the internal
capabilities required to maintain sub-microsecond or nanosecond-class
synchronization. Host limitations generally include unpredictable
software and operating system scheduling latencies that interfere with
precise clock steering, large thermal profiles that degrade standard
quartz oscillators, and a lack of specialized hardware for
bounded-latency time transfer or hardware timestamping.

The TimeCard overcomes these host limitations by completely offloading
critical timing functions---such as phase-locked loops (PLLs), holdover
tracking, and signal timestamping---to a dedicated, physically isolated
subsystem. By incorporating a TimeCard, a host system gains enhanced,
deterministic timekeeping and synchronization capabilities without
requiring a fundamental redesign of the host's primary processing
architecture.

\subsubsection{5.1.2 Implementation
Modularity}\label{implementation-modularity}

While a common physical manifestation of a TimeCard is a discrete add-in
card (such as a PCI Express card) inserted into a server chassis, a
TimeCard system is fundamentally defined by its logical interfaces and
behaviors rather than its physical constraints.

Alternate valid implementations include, but are not limited to: * A
dedicated intellectual property (IP) block directly embedded into a
System-on-Chip (SoC) or integrated onto a server motherboard. * An
external desktop or ruggedized module temporarily or permanently
connected to the host system via a hot-pluggable or peripheral interface
(e.g., USB, Thunderbolt).

Any subsystem that meets the architectural boundaries and interface
definitions defined within this standard is classified as a TimeCard for
the purposes of conformance.

\subsubsection{5.1.3 Hardware Timestamping}\label{hardware-timestamping}

To preserve the integrity and determinism of timing, it is strongly
recommended that both the providing and receiving interfaces implement
\textbf{hardware-based timestamping}. Hardware timestamping enables
timing information to be generated and measured directly within hardware
logic, avoiding non-deterministic delays caused by software stacks,
interrupt latencies, and/or operating system scheduling.

Hardware timestamping can be realized through dedicated physical
signals, such as a \textbf{Pulse-Per-Second (PPS)} output, or through
advanced in-bus implementations, such as \textbf{Precision Time
Measurement (PTM)} within modern \textbf{PCIe} architectures. By driving
the clock boundaries directly into the hardware bus, these mechanisms
enable low-latency, deterministic time delivery. This allows distributed
databases and cellular packet schedulers to achieve sub-microsecond
absolute precision, improving cross-vendor interoperability among
TimeCard and host designs.

\begin{center}\rule{0.5\linewidth}{0.5pt}\end{center}

\subsection{5.2 - Core Timing
Architecture}\label{core-timing-architecture}

At its core, every TimeCard is built around at least one
\textbf{frequency source} (with quantified stability), which serves as
the foundational source of precise timing. This source oscillator
function is complemented by one or more interfaces that enable the
TimeCard to \textbf{receive} and/or \textbf{generate} and
\textbf{distribute} time of day, phase, and frequency information to and
from the host system.

\begin{center}\rule{0.5\linewidth}{0.5pt}\end{center}

\subsection{5.3 - Inbound (Receiving) Signal
Interface}\label{inbound-receiving-signal-interface}

The \textbf{receive interface} provides a means for the TimeCard to
synchronize its oscillator function to an external reference. Depending
on the deployment environment and the required accuracy, this interface
may take multiple forms. Common examples include \textbf{Global
Navigation Satellite System (GNSS)} receivers (e.g., GPS, Galileo,
GLONASS, BeiDou), or other precision synchronization methods such as
\textbf{Precision Time Protocol (PTP)}, \textbf{Network Time Protocol
(NTP)}, \textbf{White Rabbit (WR)}, \textbf{WiWi}, \textbf{WWVB}, or
\textbf{Pulse-Per-Second (PPS)} inputs. These interfaces allow the
TimeCard to discipline its oscillator function and maintain alignment
with an external time source. The external references may or may not be
more stable or of lower noise than this oscillator function.

In some configurations, a TimeCard may operate without any inbound
external timing input. In this mode, the TimeCard functions in
\textbf{holdover}, relying solely on the stability of its internal
oscillator function to maintain accurate time over a defined interval.
Such configurations are particularly useful in environments where
external timing references are unavailable, intermittent, or
deliberately excluded for security or operational isolation to enable
redundancy.

This flexible receive architecture enables TimeCards to support a wide
use-case spectrum - from GNSS-disciplined primary time sources at the
edge of the network, to deep-indoor boundary clocks relying on PTP
*\textless\textless{} explain relevance \textgreater\textgreater**, to
autonomous isolated holdover systems - while preserving a common and
interoperable host interface standard. This ensures that a datacenter
can deploy identical host server binaries regardless of the specific
external timing source used by the TimeCard.

\begin{center}\rule{0.5\linewidth}{0.5pt}\end{center}

\subsection{5.4 - Outbound (Providing) Signal
Interface}\label{outbound-providing-signal-interface}

While the receive interface allows synchronization to an external
reference, the \textbf{providing interface} ensures that the
synchronized time and frequency are accurately distributed to the host
system.

A providing interface is a \textbf{mandatory component} of every
TimeCard. It defines the mechanism by which the TimeCard delivers time
of day, phase, frequency, or all, to the host, forming the
synchronization channel between TimeCard and host.

Depending on system requirements, the providing interface may consist of
a single interface or a combination of multiple concurrent interfaces.
Common examples include system bus standards such as \textbf{ISA},
\textbf{MCA}, \textbf{PCI}, and \textbf{PCI Express (PCIe)}, as well as
peripheral and communication interfaces such as \textbf{Serial Bus},
\textbf{USB}, \textbf{SCSI}, \textbf{PCMCIA}, or \textbf{LPT}. The
selection of interface type directly influences both the data exchange
characteristics and the precision of temporal alignment achievable by
the host.

\begin{center}\rule{0.5\linewidth}{0.5pt}\end{center}

\subsection{5.5 Management and Control Interface
(M\&CI)}\label{management-and-control-interface-mci}

In addition to the inbound and outbound signal interfaces, it is
recommended that each TimeCard include at least one \textbf{management
and control interface}. These interfaces enable configuration,
monitoring, diagnostics, firmware management, and status reporting
between the TimeCard and the host. A TimeCard without a management
interface is acceptable if its operational parameters are fixed or
pre-determined, and no runtime monitoring or control is required.

The management interface functions as the \textbf{control plane} of the
TimeCard, distinct from the \textbf{data plane} used for delivering
timing and frequency. Through this interface, the host can configure and
observe operational parameters such as oscillator state, synchronization
source selection, disciplining mode, holdover behavior, temperature
compensation, and alarm or fault conditions.

The following M\&CI busses are independent of one another, and a
TimeCard may utilize more than one kind of bus. Common examples of
management and control interfaces include but are not limited to: -
\textbf{SMBus or I²C} -- typically used for low-level configuration and
telemetry in embedded environments.\\
- \textbf{IPMI} (Intelligent Platform Management Interface) -- for
out-of-band management in server-class or rack-scale systems.\\
- \textbf{PCIe Configuration Space Registers} -- providing time-related
control and status directly over the host bus, including direct memory
access to the hardware control registers.\\
- \textbf{Serial or USB interfaces} -- enabling firmware updates,
diagnostics, or advanced telemetry access.\\
- \textbf{Network-based interfaces} such as REST, gRPC, or SNMP, for
distributed or remotely managed timing systems.

To promote interoperability and consistency, all TimeCards
\textbf{SHOULD} expose a minimum common set of registers and attributes
in a standardized format, including (but not limited to): - Current
synchronization source and state\\
- Clock disciplining status\\
- Phase and frequency offset metrics\\
- Holdover duration and expected drift\\
- Alarm and fault indicators\\
- Firmware version and build metadata

Furthermore, the management interface \textbf{SHOULD} support secure
firmware update and integrity verification mechanisms to ensure
reliability and prevent unauthorized modification. Collectively, these
management and control capabilities provide the operational transparency
and lifecycle management required for seamless integration of TimeCards
into data centers, telecom infrastructure, and AI back-end clusters. In
distributed environments, centralized management software can poll these
M\&CI endpoints to establish a global view of timing health, rapidly
identifying degraded oscillators or spoofed GNSS signals before they
impact the primary workload.

\begin{center}\rule{0.5\linewidth}{0.5pt}\end{center}

\subsection{5.6 - Power, Mechanical, and Environmental
Considerations}\label{power-mechanical-and-environmental-considerations}

The subsections apply only if the TimeCard is implemented physically,
versus for instance as a firmware function within a larger system.
\#\#\# 5.6.1 - Power Delivery - TimeCards \textbf{SHALL} define and
document input power rail voltages and tolerances (e.g., 12 V, 3.3 V).\\
- If externally powered, protection against reverse polarity
\textbf{SHALL} be included.\\
- Deterministic power-up sequencing and optional energy storage for
holdover \textbf{SHOULD} be supported.

\subsubsection{5.6.2 - Mechanical Form
Factor}\label{mechanical-form-factor}

\begin{itemize}
\tightlist
\item
  The physical envelope \textbf{SHALL} be documented.\\
\item
  Acceptable envelope forms include add-in cards
  (low-profile/full-height), mezzanine, or embedded.\\
\item
  Mounting \textbf{SHALL} withstand insertion/removal; strain relief for
  all cable ports (electrical or optical) \textbf{SHOULD} be included.\\
\item
  Faceplates \textbf{SHOULD} label at least GNSS, PPS, 10 MHz, ToD, and
  management ports and include indicator LEDs.
\end{itemize}

\subsubsection{5.6.3 - Connectors and I/O}\label{connectors-and-io}

\begin{itemize}
\tightlist
\item
  RF/timing reference signal ports \textbf{SHALL} be
  impedance-matched.\\
\item
  PPS/10 MHz electrical levels, impedance, and edge polarity (trigger on
  the rising or falling edge) \textbf{SHALL} be specified.\\
\item
  Data and management ports \textbf{SHOULD} have ESD protection and
  mechanically locking connectors.
\end{itemize}

\subsubsection{5.6.4 - Thermal Design}\label{thermal-design}

\begin{itemize}
\tightlist
\item
  Operating temperature range \textbf{SHALL} be defined; oscillator
  drift vs.~temperature \textbf{SHOULD} be characterized and
  documented.\\
\item
  Host airflow assumptions \textbf{SHOULD} be documented; heat load
  during warm-up \textbf{SHALL} be specified.\\
\item
  Over/under-temperature thresholds \textbf{SHALL} be reported via a
  management and control interface, if such an interface is provided.
\end{itemize}

\subsubsection{5.6.5 - Environmental and
Reliability}\label{environmental-and-reliability}

\begin{itemize}
\tightlist
\item
  Shock, vibration, and humidity limits \textbf{SHOULD} be stated.\\
\item
  EMC/ESD compliance \textbf{SHOULD} meet target-market standards.\\
\item
  MTBF and wear-out items \textbf{SHOULD} be documented and published.\\
\item
  Safety, labeling, and disposal requirements \textbf{SHOULD} be
  provided.
\end{itemize}

\begin{center}\rule{0.5\linewidth}{0.5pt}\end{center}

\subsection{5.7 - Reference Signals and Performance Metrics
(Normative)}\label{reference-signals-and-performance-metrics-normative}

\begin{itemize}
\tightlist
\item
  Methods defined in \textbf{Annex A (Metrics)} of this standard, which
  align with IEEE 1139, IEEE 1193, and ITU-T G.810 / G.8260,
  \textbf{SHOULD} be used for performance-metric definitions and
  analysis.
\item
  NIST Special Publication 1065 (by W.J. Riley) provides an informative
  foundation for frequency metrology.
\item
  Requirements \textbf{SHOULD} be conditioned on the physical
  characteristics of the interface type, including but not limited to:
  electrical balanced or unbalanced signaling, voltage and/or current
  thresholds, optical fiber classification, and frequency band.
\item
  ITU-T G.703 Clause 19 \textbf{MAY} be used as a reference for
  synchronous signaling.
\item
  Measurement equipment bandwidths in Hertz and trace averaging settings
  \textbf{SHALL} be explicitly documented alongside all reported
  performance results.
\end{itemize}

\subsubsection{5.7.1 - Unified Timescale
(Normative)}\label{unified-timescale-normative}

A unified timescale comes from a single oscillator function which is
published in multiple distribution formats each approximating the ideal
of the timescale to the capabilities of each format.

A TimeCard \textbf{SHALL} generate a single unified timescale and
\textbf{SHALL} publish it across all Outbound Interfaces.

All boundaries between adjacent seconds of signals from the the
providing interface of the same TimeCard instance \textbf{SHALL} align
to within a time tolerance that is documented and published.

\subsubsection{5.7.2 - Output Signal Classes
(Informative)}\label{output-signal-classes-informative}

Typical outputs include ToD, 1 PPS, 10 MHz/5 MHz, packetized time (PTP),
and host-bus time (PTM).\\
Electrical characteristics and limits \textbf{SHOULD} be published for
each.

\subsubsection{5.7.3 - Stability, Accuracy, Precision (Normative
Reporting)}\label{stability-accuracy-precision-normative-reporting}

Qualities sought include adequate stability (ADEV/TDEV/MTIE), low phase
noise, high accuracy, high precision, and fine resolution.\\
Numeric targets are intentionally unspecified; vendors \textbf{SHALL}
report measurements using: - ADEV/TDEV versus tau\\
- Time/frequency offset to reference\\
- Timestamp granularity\\
- Physical synchronization extent and conditions - And any other
measurement the manufacturer chooses to include

\subsubsection{5.7.4 - Phase Noise and Time Jitter (Normative
Reporting)}\label{phase-noise-and-time-jitter-normative-reporting}

Periodic outputs (e.g., 10 MHz) \textbf{SHOULD} include PN spectrum
vs.~offset frequency.\\
Pulse outputs (e.g., PPS), \textbf{SHALL} specify RMS and peak-to-peak
time jitter and measurement bandwidth in Hertz. These requirements
\textbf{MAY} be conditioned on intended signal and its intended use.

\subsubsection{5.7.5 - Holdover Performance
(Normative)}\label{holdover-performance-normative}

Holdover performance characterizes the stability of the TimeCard when
all primary synchronization references are disconnected or lost. -
Vendors \textbf{SHALL} publish the maximum holdover error bounds versus
elapsed time (e.g., drift in microseconds over 4, 12, and 24 hours). -
Vendors \textbf{SHALL} publish the warm-up conditions required before
the oscillator's holdover stability guarantees become valid, and the
test temperature range the holdover specification assumes.\\
- Holdover requirements strictly apply to the continuous physical drift
of the 1PPS output, assuming the TimeCard was previously locked to a
perfect, zero-noise inbound reference prior to disconnection for some
specified minimum period of time.\\
- Maximum Time Interval Error (MTIE) per ITU-T G.8260 (or G.810 App
II.5) \textbf{SHALL} be used as the definitive mathematical holdover
metric. Other auxiliary holdover metrics (such as frequency aging rate)
MAY also be measured and documented to assist system integrators.

\emph{Informative Note:} Implementers should be cautious when relying
solely on generic telecommunications boundaries (such as certain relaxed
profiles within ITU G.8262.1), as those bounds are often too loose for
the nanosecond-class strictness required by modern distributed
datacenters and AI clusters. P3335 TimeCards target significantly
tighter phase-drift boundaries.

\subsubsection{5.7.6 - Ensemble References
(Normative)}\label{ensemble-references-normative}

Implementations \textbf{SHALL} support combining multiple references
into one unified ``Ensemble'' reference.\\
Ensemble logic \textbf{MAY} provide source weights, health, and alarms
via management telemetry.

\subsubsection{5.7.7 - Large-Extent Synchronization
(Informative)}\label{large-extent-synchronization-informative}

For data-hall or campus deployments, the intent is to achieve a
specified maximum end-to-end time error while implementing calibration
as needed, and meeting cable/optical constraints as a function of the
physical dimensions (in meters) of the deployment extent. Many
independent vendors are involved, so this is a matter of overall system
design, and not of TimeCard design per se.

\subsubsection{5.7.8 - Time-Flow Narrative
(Informative)}\label{time-flow-narrative-informative}

The logical flow of time through the TimeCard architecture generally
follows these stages:

\begin{enumerate}
\def\labelenumi{\arabic{enumi}.}
\tightlist
\item
  \textbf{Ingress:} The TimeCard receives zero or more external
  references (e.g., GNSS, PTP, or PPS).
\item
  \textbf{Selection:} The system evaluates the health, stability, and
  configured priority of the incoming references and selects the most
  optimal source using a defined policy (e.g., Best Master Clock
  Algorithm).
\item
  \textbf{Disciplining (PLL):} The selected reference is fed into a
  phase-locked loop (PLL) which gracefully \textbf{\textless\textless{}
  define \textgreater\textgreater{}} disciplines the TimeCard's internal
  local oscillator. This loop filters out short-term jitter from the
  reference, relying on the high short-term stability of the local
  oscillator to provide a clean signal.
\item
  \textbf{Timescale Generation:} The disciplined oscillator drives a
  hardware counter, generating a single, unified timescale implementing
  an official timescale such as TAI or UTC.
\item
  \textbf{Egress:} The unified timescale is published across all
  Outbound Interfaces simultaneously. This includes generating physical
  PPS edges, updating memory-mapped Time of Day registers, scaling
  frequency outputs (e.g., 10 MHz), and serving host PCIe PTM
  requests---all originating deterministically from the exact same
  hardware counter.
\end{enumerate}

If the last ingress reference is lost, the PLL freezes its correction
values and the TimeCard enters \textbf{Holdover}, keeping the unified
timescale running based purely on the uncorrected drift behavior of the
local oscillator function.

\subsubsection{5.7.9 - Implementation Flexibility
(Informative)}\label{implementation-flexibility-informative}

The ``oscillator function'' need not be or contain a discrete resonator.
A nuclear frequency source may be used for applications requiring
extremenly good performasnce. A DDS or similar digital source may
suffice for cost-sensitive designs.\\
SWaP-C trade-offs are left to design and the market.

\subsubsection{5.7.10 - Conformance and Interface Definitions (Normative
Guidance)}\label{conformance-and-interface-definitions-normative-guidance}

Undefined interfaces \textbf{SHALL} normatively cite approved
\textbf{\textless\textless{} By who, and how?
\textgreater\textgreater{}} Interface Definition Documents (IDDs) for
interoperability.

Conformance testing \textbf{SHOULD} cover: - ADEV/TDEV measurement
methodology\\
- PN masks and spectra\\
- PPS alignment across all outputs\\
- PTM latency/asymmetry characterization\\
- Holdover MTIE verification\\
- Ensemble behavior under reference loss - Plus anything else the
manufacturer chooses

\begin{center}\rule{0.5\linewidth}{0.5pt}\end{center}

\subsection{5.8 - Documentation Requirements
(Normative)}\label{documentation-requirements-normative}

Manufacturers \textbf{SHALL} provide publicly available datasheets
specifying at least the following: - Which architectural principles and
constraints are implemented\\
- PLL/disciplining loop type and bandwidth\\
- All performance metrics defined in §7 (for instance stability,
accuracy, PN/jitter, holdover, ensemble)\\
- Traceability data sufficient for analysis, or an explicit
``traceability not supported'' statement\\
- All optional or conditional features provided in the implementation

\begin{center}\rule{0.5\linewidth}{0.5pt}\end{center}

\subsection{5.9 - Vendor Datasheet Checklist
(Informative)}\label{vendor-datasheet-checklist-informative}

\textbf{\textless\textless{} Explain the purpose of this section
\textgreater\textgreater{}}

{\def\LTcaptype{none} % do not increment counter
\begin{longtable}[]{@{}
  >{\raggedright\arraybackslash}p{(\linewidth - 4\tabcolsep) * \real{0.2075}}
  >{\raggedright\arraybackslash}p{(\linewidth - 4\tabcolsep) * \real{0.4528}}
  >{\raggedright\arraybackslash}p{(\linewidth - 4\tabcolsep) * \real{0.3396}}@{}}
\toprule\noalign{}
\begin{minipage}[b]{\linewidth}\raggedright
Category
\end{minipage} & \begin{minipage}[b]{\linewidth}\raggedright
Required / Recommended
\end{minipage} & \begin{minipage}[b]{\linewidth}\raggedright
Example Contents
\end{minipage} \\
\midrule\noalign{}
\endhead
\bottomrule\noalign{}
\endlastfoot
Electrical \& Power & REQUIRED & Input rails, tolerance, current draw,
sequencing \\
Mechanical & REQUIRED & Dimensions, connectors, faceplate layout \\
Thermal & REQUIRED & Operating/storage temp, airflow, warm-up load \\
Timing Performance & REQUIRED & ADEV/TDEV plots, PN spectra, holdover
MTIE \\
Synchronization Inputs & REQUIRED & GNSS/PTP/NTP/WR interface specs \\
Outputs & REQUIRED & PPS, 10 MHz, ToD specs, alignment accuracy \\
Management & RECOMMENDED & Interface type, telemetry fields, firmware
update method \\
Compliance & REQUIRED & Safety, EMC, ESD, RoHS, CE/FCC marks \\
Reliability & RECOMMENDED & MTBF, service intervals, backup power
retention \\
Documentation & REQUIRED & Quick start guide, LED legend, calibration
method \\
\end{longtable}
}

\begin{center}\rule{0.5\linewidth}{0.5pt}\end{center}

\subsection{5.10 - References (Normative)}\label{references-normative}

\begin{itemize}
\tightlist
\item
  {[}G8620{]} ITU-T G.8260 \emph{Definitions and terminology for
  synchronization in packet networks} (2015 or later).\\
\item
  {[}IEEE 1139{]} \emph{IEEE Standard Definitions of Physical Quantities
  for Fundamental Frequency and Time Metrology}.
\item
  {[}IEEE 1193{]} \emph{IEEE Guide for Measurement of Environmental
  Sensitivities of Standard Frequency Generators}
\item
  {[}MTIE{]} Stefano Bregni, \emph{Measurement of Maximum Time Interval
  Error for Telecommunications Clock Stability Characterization}, IEEE
  Transactions on Instrumentation and Measurement, Vol. 45, No.~5, Oct
  1996.\\
\item
  {[}PTPv2.1{]} IEEE 1588-2019 \emph{Precision Time Protocol (PTP)}.
\end{itemize}

\begin{center}\rule{0.5\linewidth}{0.5pt}\end{center}

\subsection{5.11 - Bibliography
(Informative)}\label{bibliography-informative}

\begin{itemize}
\tightlist
\item
  Wikipedia: {[}Precision Time Protocol{]}
  (https://en.wikipedia.org/wiki/Precision\_Time\_Protocol)\\
\item
  {[}IRS\_DID{]} DI-IPSC-81434A, \emph{Interface Requirements
  Specification Data Item Description} (1999).\\
\item
  {[}IDD\_DID{]} DI-IPSC-81436A, \emph{Interface Design Description Data
  Item Description} (1999).\\
\item
  {[}SSDD{]} DI-IPSC-81432A, \emph{System/Subsystem Design Description}
  (1999).\\
\item
  {[}SSS{]} DI-IPSC-81431A, \emph{System/Subsystem Specification}
  (2000).
\item
  NIST Special Publication 1065 by Riley
\end{itemize}

\begin{center}\rule{0.5\linewidth}{0.5pt}\end{center}

\subsection{5.12 - Notes}\label{notes}

The present P3335 standard document was initiated on 25 April 2025,
largely based on \emph{``TimeCard Architecture (Section 5) Draft
(20250424).docx''}.

\begin{center}\rule{0.5\linewidth}{0.5pt}\end{center}

\subsection{5.13 - Acronyms}\label{acronyms}

\textbf{1PPS} = One Pulse Per Second\\
\textbf{ADEV} = Allan Deviation\\
\textbf{ASIC} = Application Specific Integrated Circuit\\
\textbf{DDS} = Direct Digital Synthesis\\
\textbf{EMC} = Electro Magnetic Compatibility\\
\textbf{ESD} = Electro Static Discharge\\
\textbf{FPGA} = Field Programmable Gate Array\\
\textbf{gRPC} = Google Remote Procedure Call\\
\textbf{GNSS} = Global Navigation Satellite System\\
\textbf{I2C} = Inter-Integrated Circuit\\
\textbf{IDD} = Interface Definition Document\\
\textbf{I/O} = Input/Output\\
\textbf{IRIG} = Inter-Range Instrumentation Group\\
\textbf{ISA} = Industry Standard Architecture computer bus\\
\textbf{ITU} = International Telecommunications Union\\
\textbf{LED} = Light Emitting Diode\\
\textbf{LPT} = Line Printer Terminal\\
\textbf{M\&CI} = Management and Control Interface\\
\textbf{MCA} = Micro Channel Architecture\\
\textbf{MHz} = Megahertz (10\^{}6 Hertz)\\
\textbf{MTBF} = Mean Time Between Failure\\
\textbf{MTIE} = Maximum Time Interval Error\\
\textbf{NTP} = Network Time Protocol\\
\textbf{PCIe} = Peripheral Component Interconnect Express\\
\textbf{PCMCIA} = Personal Computer Memory Card International
Association\\
\textbf{PLL} = Phase Locked Loop\\
\textbf{PN} = Phase Noise\\
\textbf{PTM} = Precision Time Measurement (Intel)\\
\textbf{PTP} = Precision Time Protocol\\
\textbf{REST} = Representational State Transfer\\
\textbf{RMS} = Root Mean Square\\
\textbf{SCSI} = Small Computer System Interface\\
\textbf{SMB} = System Management Bus\\
\textbf{SNMP} = Simple Network Management Protocol\\
\textbf{SoC} = System on a Chip\\
\textbf{SWaP-C} = Size, Weight, Power, and Cost\\
\textbf{TAI} = International Atomic Time (in French)\\
\textbf{TDEV} = Time Deviation\\
\textbf{ToD} = Time of Day\\
\textbf{USB} = Universal Serial Bus\\
\textbf{UTC} = Coordinated Universal Time\\
\textbf{WG} = Working Group\\
\textbf{WiWi} = Wireless two-Way interferometry\\
\textbf{WR} = White Rabbit\\
\textbf{WWVB} = Radio Station WWVB

\begin{center}\rule{0.5\linewidth}{0.5pt}\end{center}

\textbf{End of Document}

\section{Performance Metrics
(Normative)}\label{performance-metrics-normative}

This clause defines the performance reporting and characterization
requirements for TimeCard implementations. The intent is to enable
measurable, transparent, and comparable timing behavior across vendors.
Performance metrics shall be reported in vendor documentation.

\begin{center}\rule{0.5\linewidth}{0.5pt}\end{center}

\subsection{6.1 Overview}\label{overview}

TimeCard performance must quantify the precision, stability, and
accuracy with which the device maintains and delivers time and/or
frequency in a context aligned with the performance requirements of the
intended application.

Performance metrics for the oscillator(s) utilized within the TimeCard
shall be provided. These shall include frequency accuracy and
temperature stability.

Performance metrics applicable for the characterization of the operation
within the applicable use case should be provided. Those metrics should
be characterized under both locked (synchronized/syntonized) and
holdover operating conditions.

Manufacturers may report the following metrics: - Frequency stability
(ADEV, TDEV) - Maximum Time Interval Error (MTIE) - Phase noise (PN) -
Jitter and alignment (PPS and ToD) - Accuracy and drift relative to
reference - Holdover behavior - Environmental sensitivity (to
temperature, voltage, vibration)

All measurements shall be made by testing equipment traceable to the
relevant standards, such as UTC(NIST), and conform to relevant
methodologies e.g.~ITU-T G.810/G.8260.

\subsection{6.2 Benchmarking and Testing Oscillators on the
TimeCard}\label{benchmarking-and-testing-oscillators-on-the-timecard}

To validate the capabilities of the onboard oscillator(s) within the
context of the TimeCard system, comprehensive testing procedures shall
be conducted. The presence of the TimeCard inside a host system (such as
a PCIe slot in a standard server) introduces environmental stressors
such as thermal gradients, voltage fluctuations, and vibrations, which
must be accounted for during benchmarking.

\subsubsection{6.2.1 Test Environment and
Setup}\label{test-environment-and-setup}

Benchmarking shall be performed using industry-standard precision
measurement equipment to ensure reproducibility and traceability: -
\textbf{Reference Sources:} An atomic reference standard (e.g., Cesium
or Rubidium clock) or a highly stable GNSS-disciplined oscillator
(GNSSDO) tracing back to UTC. - \textbf{Measurement Instruments:}
High-resolution Time Interval Counters (TIC), Phase Noise Analyzers, and
Oscillators with low intrinsic noise. - \textbf{Environmental Controls:}
Thermal chambers and vibration tables for simulating realistic operation
conditions.

\subsubsection{6.2.2 Core Performance
Characterization}\label{core-performance-characterization}

The following tests are standard for evaluating the intrinsic
performance of the oscillator on the TimeCard:

\begin{enumerate}
\def\labelenumi{\arabic{enumi}.}
\tightlist
\item
  \textbf{Fractional Frequency Stability (ADEV, TDEV):}

  \begin{itemize}
  \tightlist
  \item
    Measured by comparing the TimeCard's oscillator output (e.g., 10 MHz
    or 1 PPS) against the reference source over various observation
    intervals ((\tau)).
  \item
    Allan Deviation (ADEV) and Time Deviation (TDEV) metrics shall be
    plotted to characterize short-term, medium-term, and long-term
    instability (white noise, flicker noise, and random walk).
  \end{itemize}
\item
  \textbf{Phase Noise Testing:}

  \begin{itemize}
  \tightlist
  \item
    Evaluated in the frequency domain using a Phase Noise Analyzer to
    quantify short-term phase fluctuations.
  \item
    Typically reported at standard offset frequencies (e.g., 1 Hz, 10
    Hz, 100 Hz, 1 kHz, 10 kHz, 100 kHz) from the carrier.
  \end{itemize}
\item
  \textbf{Maximum Time Interval Error (MTIE) and Time Variance (TVAR):}

  \begin{itemize}
  \tightlist
  \item
    Assesses the peak-to-peak variation in time delay of the
    oscillator's output relative to an ideal reference over a specified
    observation window.
  \item
    Crucial for determining whether the oscillator's wander meets
    telecom and networking standards (e.g., ITU-T G.8262).
  \end{itemize}
\end{enumerate}

\subsubsection{6.2.3 Operational and Holdover
Benchmarking}\label{operational-and-holdover-benchmarking}

The interaction between the synchronization subsystem and the oscillator
must be tested under dynamic conditions.

\begin{enumerate}
\def\labelenumi{\arabic{enumi}.}
\tightlist
\item
  \textbf{Lock Acquisition and Tracking:}

  \begin{itemize}
  \tightlist
  \item
    Measure the time required to achieve syntonization and
    synchronization from a cold start and a warm start.
  \item
    Characterize phase transients and frequency overshoot during the
    locking phase.
  \end{itemize}
\item
  \textbf{Holdover Performance:}

  \begin{itemize}
  \tightlist
  \item
    \textbf{Procedure:} Allow the TimeCard to reach a fully stabilized
    locked state using a GNSS or PTP reference. Disconnect the reference
    source to force the system into holdover mode.
  \item
    \textbf{Measurement:} Record the time drift (phase error)
    accumulation relative to the reference over specific periods (e.g.,
    1 hour, 4 hours, 24 hours).
  \item
    The reported holdover class must specify the maximum time error
    observed over a defined duration and temperature profile.
  \end{itemize}
\end{enumerate}

\subsubsection{6.2.4 Environmental Stress
Testing}\label{environmental-stress-testing}

Since TimeCards are intended for deployment in standard computing
hardware, the oscillators must be benchmarked under simulated data
center or edge environments.

\begin{enumerate}
\def\labelenumi{\arabic{enumi}.}
\tightlist
\item
  \textbf{Thermal Stability Profile:}

  \begin{itemize}
  \tightlist
  \item
    Expose the TimeCard to a defined temperature ramp (e.g., -5°C to
    +55°C depending on chassis requirements).
  \item
    For high-precision applications, determine the dynamic thermal
    gradient response (e.g., drift resulting from a 1°C/minute change)
    and steady-state frequency over temperature (Δf/f vs.~ΔT).
  \end{itemize}
\item
  \textbf{Power Supply and Voltage Variation:}

  \begin{itemize}
  \tightlist
  \item
    The PCIe bus provides power, but servers can exhibit significant
    voltage ripple.
  \item
    Test frequency stability and phase noise under ±5\% or ±10\% input
    voltage fluctuations, verifying the effectiveness of onboard Low
    Dropout Regulators (LDOs) and power-filtering circuitry.
  \end{itemize}
\item
  \textbf{Vibration and Airflow:}

  \begin{itemize}
  \tightlist
  \item
    Data center chassis fans cause mechanical vibration and rapid,
    localized temperature shifts due to airflow.
  \item
    Benchmark the oscillator using a shaker table matching typical
    server rack vibration profiles (e.g., ETSI EN 300 019).
  \item
    Measure vibration-induced phase noise and dynamic frequency shifts.
  \end{itemize}
\end{enumerate}

\subsubsection{6.2.5 PCIe and System Integration
Impact}\label{pcie-and-system-integration-impact}

It is critical to test the oscillator's performance when integrated
within a standard PCIe environment, rather than isolated on a test
bench: - Compare baseline (standalone test bench) oscillator performance
to performance while executing heavy PCIe data transfers. - Verify that
standard clock domains within the server do not induce crosstalk or
correlated jitter on the TimeCard's precision timing signals.

\section{7. Timing Interfaces
(Normative)}\label{timing-interfaces-normative}

This chapter defines the physical, electrical, and logical interfaces
through which a TimeCard exchanges timing and frequency information with
external systems. These interfaces form the foundation of
interoperability between diverse TimeCards, host platforms, and external
timing networks. All implementations SHALL adhere to the standardized
behaviors defined herein to facilitate cross-vendor compatibility and
predictable operation.

\subsection{7.1 Overview}\label{overview-1}

A TimeCard typically incorporates three fundamental classes of
interfaces: 1. \textbf{Receive (Inbound) Interfaces:} Dedicated to
acquiring timing signals from external authoritative references. 2.
\textbf{Providing (Outbound) Interfaces:} Dedicated to distributing
disciplined time, phase, and frequency representations to the host
system and other downstream consumers. 3. \textbf{Management and Control
Interfaces:} Non-timing-critical data paths used for configuration,
telemetry, security, and diagnostics.

Each class of interface possesses distinct electrical, signaling,
protocol, and performance requirements, which are detailed in the
following subclauses.

\subsection{7.2 Receive (Inbound)
Interfaces}\label{receive-inbound-interfaces}

The receive interface enables the TimeCard to lock its internal
oscillator function to one or more external references, establishing
traceability and synchronization to a higher-order time source.

\subsubsection{7.2.1 Supported Reference
Types}\label{supported-reference-types}

The following external references are recognized for TimeCard
synchronization. A TimeCard MAY support one or more of these inputs: -
\textbf{GNSS Receivers:} (e.g., GPS, Galileo, GLONASS, BeiDou) supplying
1 Pulse Per Second (1PPS) and Time of Day (ToD) streams such as NMEA or
binary protocols. - \textbf{Precision Time Protocol (PTP):} Packet-based
synchronization via IEEE Std 1588 over Ethernet or PCIe networks. -
\textbf{Network Time Protocol (NTP):} Packet-based synchronization for
administrative or lower-accuracy time distribution via RFC 5905. -
\textbf{White Rabbit (WR):} Sub-nanosecond synchronization via
deterministic Ethernet variants. - \textbf{WiWi:} Wireless precision
timing leveraging IEEE 802.1AS extensions. - \textbf{Regional Radio
Standards:} (e.g., WWVB, DCF77, MSF) for localized low-frequency time
distribution. - \textbf{Direct Electrical References:} 1PPS and
continuous frequency (e.g., 10 MHz) continuous-wave or square-wave
inputs for direct discrete phase locking or localized chaining.

\subsubsection{7.2.2 Interface
Characteristics}\label{interface-characteristics}

For each physical receive interface implemented, the manufacturer SHALL
rigorously specify the following characteristics in the product
documentation: - \textbf{Electrical level:} (e.g., 3.3V CMOS, LVTTL,
RS-422, sinewave dBm). - \textbf{Input impedance and line termination:}
(e.g., 50 Ω, 75 Ω, high-Z). - \textbf{Signal polarity and edge
definition:} For pulse inputs, defining whether the rising or falling
edge acts as the on-time phase marker.

Additionally, the lock acquisition time parameters and holdover
entry/exit behaviors are configurable properties of the TimeCard and
SHALL be documented.

\subsubsection{7.2.3 Input Redundancy and
Selection}\label{input-redundancy-and-selection}

TimeCards supporting multiple reference inputs SHALL implement source
selection and failover logic (e.g., an implementation of the Best Master
Clock Algorithm or a proprietary multi-source weighting ensemble).

The source priority, weighting, and failover policy SHALL be
configurable via the management interface. During a reference transition
(failover), the TimeCard internal oscillator SHOULD maintain phase
continuity; however, derivative outputs are permitted minor bounded
discontinuities as defined by the manufacturer's specification.

\subsubsection{7.2.4 Sensitivity and
Thresholds}\label{sensitivity-and-thresholds}

The physical receive path SHALL tolerate nominal input amplitude
variations of at least ±20\% and reject high-frequency noise consistent
with ITU-T G.8260 sensitivity definitions.

Loss-of-signal (LOS) detection boundaries and subsequent signal
reacquisition thresholds SHALL be explicitly documented in the vendor's
datasheet.

\subsection{7.3 Providing (Outbound)
Interfaces}\label{providing-outbound-interfaces}

The providing interface is responsible for delivering the unified,
synchronized timescale generated by the TimeCard to the host platform or
external physical devices.

\subsubsection{7.3.1 Output Signal Classes}\label{output-signal-classes}

Typical providing interfaces include, but are not limited to: -
\textbf{Pulse-Per-Second (PPS):} The primary second-aligned discrete
reference pulse. - \textbf{Frequency Outputs:} Continuous clocks (e.g.,
10 MHz, 5 MHz) used for frequency synchronization. - \textbf{Time of Day
(ToD):} Serial-formatted absolute time messages (e.g., NMEA, IRIG-B, or
proprietary binary telemetry). - \textbf{Packetized Time (PTP):} IEEE
1588 traffic delivered via Ethernet ports or directly into the host via
PCIe networking abstractions. - \textbf{In-Bus Timing (PTM):} PCIe
Precision Time Measurement transactions directly delivered to the host
bridging architecture. - \textbf{Synchronous Clocks:} Discrete
phase-aligned clocks for driving localized FPGA or SoC logic domains.

\subsubsection{7.3.2 Electrical and Protocol
Standards}\label{electrical-and-protocol-standards}

All physical outputs SHALL comply with the following baseline signaling
levels unless otherwise documented for specialized physical
environments: - \textbf{PPS and frequency outputs:} 3.3 V CMOS or LVTTL
utilizing 50 Ω source impedance. - \textbf{ToD / Serial outputs:}
RS-232, RS-422, or RS-485 compliant signaling levels. - \textbf{PCIe /
PTM interfaces:} Compliant with PCI-SIG 5.0 specifications (or later). -
\textbf{Ethernet / PTP interfaces:} Compliant with IEEE Std 1588-2019
and/or IEEE Std 802.1AS-2020.

For each implemented providing interface, the manufacturer SHALL
document the signal rise/fall times, maximum intrinsic jitter, source
impedance, and absolute drive capability.

\subsubsection{7.3.3 Synchronization
Accuracy}\label{synchronization-accuracy}

All providing output signals SHALL align coherently to the single
unified timescale generated natively by the TimeCard as described in the
Architecture clauses.

For compliant high-precision designs, the 1PPS output alignment error
SHALL NOT exceed 1 nanosecond RMS relative to the perfectly aligned
internal master reference point. Output clocks SHOULD exhibit
deterministic phase behavior and strictly bounded frequency deviations
during any internal source failover event.

\subsubsection{7.3.4 Hardware
Timestamping}\label{hardware-timestamping-1}

Providing interfaces SHALL support hardware-assisted timestamping where
technically feasible within the protocol, including: - Direct physical
PPS edge capture. - In-band PCIe PTM transaction stamping. - PTP
transmit (egress) and receive (ingress) hardware packet timestamps.

Timestamping accuracy validation SHALL be performed using measurement
equipment directly traceable to at least one recognized National
Metrology Institute (e.g., NIST, NPL, or PTB).

\subsection{7.4 Management and Control
Interfaces}\label{management-and-control-interfaces}

Management interfaces define the logical communications paths dedicated
to device configuration, telemetry polling, and firmware lifecycle
control. While these interfaces do not carry the critical
synchronization timepath, they are indispensable for operational
deployment.

\subsubsection{7.4.1 Common Management
Interfaces}\label{common-management-interfaces}

{\def\LTcaptype{none} % do not increment counter
\begin{longtable}[]{@{}
  >{\raggedright\arraybackslash}p{(\linewidth - 6\tabcolsep) * \real{0.2143}}
  >{\raggedright\arraybackslash}p{(\linewidth - 6\tabcolsep) * \real{0.2500}}
  >{\raggedright\arraybackslash}p{(\linewidth - 6\tabcolsep) * \real{0.2143}}
  >{\raggedright\arraybackslash}p{(\linewidth - 6\tabcolsep) * \real{0.3214}}@{}}
\toprule\noalign{}
\begin{minipage}[b]{\linewidth}\raggedright
Interface
\end{minipage} & \begin{minipage}[b]{\linewidth}\raggedright
Typical Use
\end{minipage} & \begin{minipage}[b]{\linewidth}\raggedright
Relative Bandwidth
\end{minipage} & \begin{minipage}[b]{\linewidth}\raggedright
Example Use Case
\end{minipage} \\
\midrule\noalign{}
\endhead
\bottomrule\noalign{}
\endlastfoot
\textbf{SMBus / I²C} & Low-level baseboard configuration & Low & Host
BMC telemetry, power sequencing, clock control \\
\textbf{IPMI / NC-SI} & Out-of-band server management & Medium &
Platform health monitoring, remote chassis querying \\
\textbf{PCIe Config Space} & In-band direct host control & High &
Runtime MMIO status registers, PCIe PTM enablement \\
\textbf{Serial / USB} & Maintenance and physical access & Medium &
Localized firmware flashing, debug CLI access \\
\textbf{REST / gRPC / SNMP} & Remote network orchestration & Variable &
Distributed datacenter timing management systems \\
\end{longtable}
}

\subsubsection{7.4.2 Management
Capabilities}\label{management-capabilities}

Where a management interface is provided, it SHALL allow read and write
access to at minimum the following operational vectors: - Active
synchronization source selection and current fallback state. - Internal
oscillator disciplining mode (e.g., locked, holdover, free-run, warming
up). - Phase and frequency offset statistics relative to the reference.
- Active alarm registers, event logs, and threshold warnings. - Firmware
versioning, hardware revisions, and active security status.

\subsubsection{7.4.3 Security and Sanitization
Requirements}\label{security-and-sanitization-requirements}

Given the critical nature of timing infrastructure, management
interfaces SHALL adhere to strict security protocols: - All remote or
networked management channels SHALL require standard authentication
mechanisms (e.g., username/password, cryptographic tokens, or mutual TLS
certificates). - Firmware payload updates SHALL mandate cryptographic
signing and boot-time integrity verification to prevent unauthorized
code execution. - Telemetry data crossing shared or external network
boundaries SHOULD be encrypted. - \textbf{Data Sanitization:} TimeCards
SHOULD support a cryptographic erase or sanitization procedure. Upon
receiving a specific management command, or upon detection of a physical
tamper event, the system SHALL permanently erase all sensitive internal
state data, private keys, and operational logs, satisfying requirements
for environments where the overall system must forget all internal data
upon a sanitization procedure.

\subsection{7.5 Physical and Logical
Integration}\label{physical-and-logical-integration}

\subsubsection{7.5.1 Connectors and
Labeling}\label{connectors-and-labeling}

Physical TimeCard front panels (brackets) SHALL label all timing and
management ports clearly. For example:
\texttt{{[}GNSS{]}\ {[}PPS\ OUT{]}\ {[}10MHz{]}\ {[}ToD{]}\ {[}MGMT{]}\ {[}USB{]}}

Manufacturers SHOULD adopt standardized color coding and iconography in
accordance with Open Compute Project Time Appliance Project (OCP-TAP)
guidelines where applicable.

\subsubsection{7.5.2 Interoperability Goals and
Requirements}\label{interoperability-goals-and-requirements}

\begin{itemize}
\tightlist
\item
  All providing interfaces on a single TimeCard SHALL publish the same
  mathematically aligned, unified timescale.
\item
  Interfaces SHOULD support broad cross-vendor interoperability by
  utilizing only the common electrical and logical formats specified
  herein.
\item
  \textbf{Architectural Objective:} While vendors MAY implement
  proprietary or optional extensions to the TimeCard interface, it is a
  primary design objective of this standard that such extensions do not
  natively disrupt or degrade baseline interoperability with standard
  host systems.
\end{itemize}

\subsection{7.6 Summary}\label{summary}

Timing interfaces construct the foundational boundaries of TimeCard
interoperability, facilitating precise, deterministic, and
hardware-agnostic synchronization. Compliance with the physical and
logical interface definitions in this chapter realizes the goal that any
TimeCard---regardless of underlying component design or vendor
origin---can seamlessly integrate into diverse host environments,
sustain traceable absolute time accuracy, and distribute stable
frequency and phase alignments.

\begin{center}\rule{0.5\linewidth}{0.5pt}\end{center}

\section{8. Control Interfaces
(Normative)}\label{control-interfaces-normative}

This chapter defines the physical, logical, and protocol-level
\textbf{control interfaces} used to configure, monitor, and manage
TimeCard devices.

\begin{center}\rule{0.5\linewidth}{0.5pt}\end{center}

\subsection{8.1 Overview}\label{overview-2}

Control interfaces form the out-of-band and in-band \textbf{management
plane} of the TimeCard architecture. Unlike the timing interfaces
defined in Clause 7, which operate strictly within the data plane to
deliver phase and frequency synchronization, control interfaces are
responsible for: - Configuration of synchronization parameters and state
machines. - Status monitoring, health telemetry, and fault reporting. -
Firmware provisioning and lifecycle security management. - Traceability
reporting and calibration operations.

Every compliant TimeCard implementation SHALL provide at least one
accessible control interface. This interface MAY be exposed through a
physical baseboard connector, a host-exposed bus interface, or a
network-based protocol suite.

\begin{center}\rule{0.5\linewidth}{0.5pt}\end{center}

\subsection{8.2 Control Interface
Classes}\label{control-interface-classes}

Control interfaces are categorized based on their communication topology
and data granularity.

{\def\LTcaptype{none} % do not increment counter
\begin{longtable}[]{@{}
  >{\raggedright\arraybackslash}p{(\linewidth - 4\tabcolsep) * \real{0.2222}}
  >{\raggedright\arraybackslash}p{(\linewidth - 4\tabcolsep) * \real{0.3889}}
  >{\raggedright\arraybackslash}p{(\linewidth - 4\tabcolsep) * \real{0.3889}}@{}}
\toprule\noalign{}
\begin{minipage}[b]{\linewidth}\raggedright
Class
\end{minipage} & \begin{minipage}[b]{\linewidth}\raggedright
Description
\end{minipage} & \begin{minipage}[b]{\linewidth}\raggedright
Typical Use
\end{minipage} \\
\midrule\noalign{}
\endhead
\bottomrule\noalign{}
\endlastfoot
\textbf{Local Hardware Control} & Direct access via SMBus, I²C, I3C, or
GPIO. & Baseboard-level configuration, power sequencing, or low-level
telemetry by a BMC. \\
\textbf{Host-Integrated Control} & In-band access via PCIe MMIO or USB.
& Runtime synchronization control by host OS software or drivers. \\
\textbf{Out-of-Band (OOB) Control} & NC-SI, IPMI, or dedicated serial
interface. & Independent platform management regardless of the host OS
state. \\
\textbf{Remote Network Control} & REST, gRPC, Redfish, or SNMP. &
Distributed system orchestration or cloud-scale management. \\
\end{longtable}
}

A TimeCard MAY implement one or more of these control classes
simultaneously to support sophisticated datacenter abstraction models.

\begin{center}\rule{0.5\linewidth}{0.5pt}\end{center}

\subsection{8.3 Minimum Functional
Requirements}\label{minimum-functional-requirements}

Where a control interface is implemented, it SHALL expose at least the
following minimum common set of capabilities, restricted by the
bandwidth and constraints of the specific medium:

{\def\LTcaptype{none} % do not increment counter
\begin{longtable}[]{@{}
  >{\raggedright\arraybackslash}p{(\linewidth - 2\tabcolsep) * \real{0.4400}}
  >{\raggedright\arraybackslash}p{(\linewidth - 2\tabcolsep) * \real{0.5600}}@{}}
\toprule\noalign{}
\begin{minipage}[b]{\linewidth}\raggedright
Function
\end{minipage} & \begin{minipage}[b]{\linewidth}\raggedright
Description
\end{minipage} \\
\midrule\noalign{}
\endhead
\bottomrule\noalign{}
\endlastfoot
\textbf{Source Selection} & Select, prioritize, or configure the active
synchronization reference (e.g., GNSS, PTP, PPS). \\
\textbf{Mode Control} & Read and modify the current operational state
machine (e.g., Locked, Holdover, Free-running, Warm-up). \\
\textbf{Clock Discipline Settings} & Adjust PLL or disciplining loop
time constants and bandwidth parameters. \\
\textbf{Telemetry Access} & Retrieve metrics such as fractional
frequency offset, internal temperatures, instantaneous phase error, and
stability statistics. \\
\textbf{Alarm and Event Logs} & Report critical operational events such
as loss of reference, spoofing detection, or firmware integrity
faults. \\
\textbf{Firmware Management} & Support cryptographic firmware update
triggers, version reporting, and rollback. \\
\end{longtable}
}

\begin{center}\rule{0.5\linewidth}{0.5pt}\end{center}

\subsection{8.4 Communication Protocols}\label{communication-protocols}

\subsubsection{8.4.1 SMBus / I²C / I3C}\label{smbus-iuxb2c-i3c}

\begin{itemize}
\tightlist
\item
  Provides low-speed, low-latency configuration and sensor access,
  typically used by a Baseboard Management Controller (BMC).
\item
  The address space and register map SHOULD be modeled consistently
  across compliant implementations to ease integration.
\item
  Bus arbitration, clock stretching, and timing compliance SHALL follow
  the respective SMBus or MIPI I3C specifications.
\end{itemize}

\subsubsection{8.4.2 PCIe Configuration and
MMIO}\label{pcie-configuration-and-mmio}

\begin{itemize}
\tightlist
\item
  Allows in-band, high-speed access to TimeCard control registers
  directly from the host CPU.
\item
  The register layout SHOULD comply with recognized industry standard
  memory maps (such as the OCP-TAP Control Register Map) where
  applicable.
\item
  If the TimeCard supports PCIe Precision Time Measurement (PTM), the
  control interface SHALL present the necessary standard capability
  structures to the host to enable PTM negotiation.
\end{itemize}

\subsubsection{8.4.3 IPMI / NC-SI}\label{ipmi-nc-si}

\begin{itemize}
\tightlist
\item
  Enables robust out-of-band control independent of the host processor
  or operating system.
\item
  TimeCard devices establishing IPMI communication SHOULD implement an
  extended command set for time synchronization management (e.g.,
  querying clock state, initiating reference failover).
\item
  IPMI commands MAY support authenticated sessions utilizing pre-shared
  platform credentials.
\end{itemize}

\subsubsection{8.4.4 REST / gRPC / SNMP /
Redfish}\label{rest-grpc-snmp-redfish}

\begin{itemize}
\tightlist
\item
  Used for high-level, network-based management and large-scale
  telemetry aggregation.
\item
  REST, Redfish, and gRPC APIs SHOULD conform to established Open
  Management Schemas (OMS).
\item
  SNMP implementations SHOULD expose standard OIDs for timing health,
  synchronization source, and hardware/firmware versions.
\end{itemize}

\subsubsection{8.4.5 Serial and USB}\label{serial-and-usb}

\begin{itemize}
\tightlist
\item
  Serial (RS-232/RS-422/RS-485) interfaces SHOULD provide a
  human-readable command-line shell or SCPI-like syntax for localized
  debugging.
\item
  USB connections MAY present as standardized Communications Device
  Class (CDC) / virtual COM ports to facilitate maintenance or localized
  firmware provisioning.
\end{itemize}

\begin{center}\rule{0.5\linewidth}{0.5pt}\end{center}

\subsection{8.5 Control Register
Structure}\label{control-register-structure}

The logical control registers listed below represent the fundamental
baseline for low-level interaction between the host/BMC and the
TimeCard. While physical implementations differ (e.g., I2C addresses
vs.~PCIe MMIO offsets), the logical concepts SHALL be present:

{\def\LTcaptype{none} % do not increment counter
\begin{longtable}[]{@{}
  >{\raggedright\arraybackslash}p{(\linewidth - 6\tabcolsep) * \real{0.2200}}
  >{\raggedright\arraybackslash}p{(\linewidth - 6\tabcolsep) * \real{0.2800}}
  >{\raggedright\arraybackslash}p{(\linewidth - 6\tabcolsep) * \real{0.1800}}
  >{\raggedright\arraybackslash}p{(\linewidth - 6\tabcolsep) * \real{0.3200}}@{}}
\toprule\noalign{}
\begin{minipage}[b]{\linewidth}\raggedright
Logical Register Concept
\end{minipage} & \begin{minipage}[b]{\linewidth}\raggedright
Description
\end{minipage} & \begin{minipage}[b]{\linewidth}\raggedright
Access
\end{minipage} & \begin{minipage}[b]{\linewidth}\raggedright
Example Units
\end{minipage} \\
\midrule\noalign{}
\endhead
\bottomrule\noalign{}
\endlastfoot
\texttt{SYNC\_SRC} & Active external synchronization source ID & R/W &
Enumeration \\
\texttt{PLL\_STATE} & Current disciplining state (Locked, Holdover,
Free) & R & Enumeration \\
\texttt{PHASE\_ERR} & Instantaneous phase error relative to the primary
reference & R & nanoseconds \\
\texttt{TEMP} & Active local oscillator temperature & R & °C \\
\texttt{MTIE\_EST} & Current estimated holdover drift/MTIE boundary & R
& nanoseconds \\
\texttt{FIRMWARE\_VER} & Active running firmware version string & R &
ASCII / Hex \\
\texttt{UPDATE\_CMD} & Trigger flag for firmware update sequences & W &
Boolean \\
\texttt{SECURE\_MODE} & Runtime security lockdown indicator & R/W &
Boolean \\
\end{longtable}
}

All registers SHOULD adhere to consistent endianness and word alignment
across a vendor's product line.

\begin{center}\rule{0.5\linewidth}{0.5pt}\end{center}

\subsection{8.6 Security and Access
Control}\label{security-and-access-control}

\emph{Note: The requirements in this subclause apply conditionally. They
are REQUIRED only for TimeCards marketed or designated for secure
infrastructure deployments. For isolated or lab-grade physical
environments, these features MAY be omitted to reduce complexity.}

Where security hardening is implemented, control interfaces operate as
primary attack surfaces and require robust cryptographic protection.

\subsubsection{8.6.1 Authentication and
Authorization}\label{authentication-and-authorization}

\begin{itemize}
\tightlist
\item
  Networked control protocols SHALL support user or machine
  authentication.
\item
  Privilege levels SHOULD distinguish between read-only monitoring
  access, operator configuration access, and administrator provisioning
  access.
\item
  Access credentials and cryptographic keys SHOULD be stored securely
  within a hardware-backed enclave or discrete TPM.
\end{itemize}

\subsubsection{8.6.2 Secure Firmware and
Configuration}\label{secure-firmware-and-configuration}

\begin{itemize}
\tightlist
\item
  Firmware payload images SHALL be digitally signed using verifiable
  certificates.
\item
  A secure boot sequence SHOULD verify firmware integrity (cryptographic
  hash matches) at initialization before granting control to the
  TimeCard operating logic.
\end{itemize}

\subsubsection{8.6.3 Network Security}\label{network-security}

\begin{itemize}
\tightlist
\item
  Remote IP-based interfaces (REST, Redfish, gRPC) SHOULD utilize
  encrypted transport layers (e.g., TLS).
\item
  Replay and injection protection SHOULD be implemented for operations
  traversing shared networking environments.
\end{itemize}

\begin{center}\rule{0.5\linewidth}{0.5pt}\end{center}

\subsection{8.7 Event and Telemetry
Management}\label{event-and-telemetry-management}

To facilitate continuous performance observability, the control
interface SHOULD support structured event and telemetry generation.

\subsubsection{8.7.1 Event Reporting}\label{event-reporting}

\begin{itemize}
\tightlist
\item
  Generated operational events SHALL include a local timestamp, a
  severity classification, and a source identifier.
\item
  Critical events SHOULD include, but are not limited to:

  \begin{itemize}
  \tightlist
  \item
    Reference loss (LOS)
  \item
    Phase lock achieved or lost
  \item
    Holdover entry or exit
  \item
    Temperature threshold alarms
  \item
    Firmware update success or failure.
  \end{itemize}
\end{itemize}

\subsubsection{8.7.2 Telemetry Streaming}\label{telemetry-streaming}

\begin{itemize}
\tightlist
\item
  Continuous telemetry MAY be streamed via IPMI, REST, or gRPC for
  integration into centralized infrastructure dashboards.
\item
  Telemetry generation rates SHALL be configurable by the administrator.
\item
  Time synchronization statistics (such as ADEV or instantaneous phase
  offsets) SHOULD be strictly timestamped and aligned to the TimeCard's
  unified timescale.
\end{itemize}

\subsubsection{8.7.3 Logging and
Traceability}\label{logging-and-traceability}

\begin{itemize}
\tightlist
\item
  Critical alarm logs SHOULD be preserved in non-volatile memory across
  power cycles.
\item
  Exported log files MAY include cryptographic signatures to validate
  log integrity during forensic analysis.
\item
  Traceability metadata (such as calibration constants and firmware
  build hashes) SHOULD be accessible to support operational compliance
  audits.
\end{itemize}

\begin{center}\rule{0.5\linewidth}{0.5pt}\end{center}

\subsection{8.8 Interoperability and
Extensions}\label{interoperability-and-extensions}

All implemented control interfaces SHALL comply with the baseline schema
defined by the overarching IEEE P3335 architecture.

Manufacturers MAY implement proprietary extensions, custom registers, or
vendor-specific telemetry strings. However, as an architectural
objective, these optional extensions SHOULD NOT compromise or interfere
with the baseline interoperability of the standardized registers and
capabilities.

\begin{center}\rule{0.5\linewidth}{0.5pt}\end{center}

\subsection{8.9 Summary}\label{summary-1}

Control interfaces construct the necessary framework for secure,
consistent, and vendor-neutral lifecycle management of TimeCard systems.

By standardizing configuration and telemetry logic, this framework
facilitates the integration of diverse TimeCards into heterogeneous
server fleets and telecommunications infrastructures. This architectural
consistency promotes long-term maintainability, traceability, and
operational reliability throughout the lifespan of precision time
synchronization networks.

\begin{center}\rule{0.5\linewidth}{0.5pt}\end{center}

\textbf{End of Chapter -- Control Interfaces}

\section{9. Environmental Specifications
(Normative)}\label{environmental-specifications-normative}

This chapter defines the environmental, mechanical, and reliability
parameters applicable to compliant TimeCard implementations. The
objective is to establish predictable hardware behavior under diverse
deployment conditions---ranging from controlled data centers and
telecommunications facilities to industrial or field
installations---while maintaining timing precision and conformance to
the unified timescale.

\begin{center}\rule{0.5\linewidth}{0.5pt}\end{center}

\subsection{9.1 Overview}\label{overview-3}

TimeCards are precision timing subsystems whose performance is directly
influenced by environmental factors such as ambient temperature, thermal
gradients, humidity, mechanical shock, vibration, and electromagnetic
interference.

All implementations SHALL specify their environmental operating limits,
test methodologies, and mitigation measures in compliance with the
standards to which the manufacturer claims conformity.

Environmental performance SHOULD be validated through qualification
testing to verify that frequency, phase, and absolute time accuracy
remain stable and within stated bounds over the entire declared
operating profile.

\begin{center}\rule{0.5\linewidth}{0.5pt}\end{center}

\subsection{9.2 Operating Conditions}\label{operating-conditions}

\subsubsection{9.2.1 Temperature}\label{temperature}

\begin{itemize}
\tightlist
\item
  \textbf{Operating Range:} Each TimeCard SHALL explicitly define a
  continuous operating temperature range suited for its target
  deployment (e.g., 0°C to +55°C for data center use, or −40°C to +85°C
  for industrial use).\\
\item
  \textbf{Storage Range:} Vendors SHALL specify safe, non-operational
  storage temperatures and humidity levels.\\
\item
  \textbf{Thermal Stability:} Local oscillator frequency drift versus
  temperature (temperature coefficient) SHALL be characterized and
  documented.\\
\item
  \textbf{Compensation:} Precision devices employing Oven-Controlled
  Crystal Oscillators (OCXO) or Temperature-Compensated Crystal
  Oscillators (TCXO) SHOULD document their integrated temperature
  compensation logic or calibration tables.\\
\item
  \textbf{Thermal Alarms:} Active management telemetry, if present,
  SHALL include configurable over-temperature and under-temperature
  warning thresholds.
\end{itemize}

\subsubsection{9.2.2 Humidity}\label{humidity}

\begin{itemize}
\tightlist
\item
  \textbf{Operating Humidity:} Devices SHALL be specified to operate
  safely between 5\% to 95\% non-condensing relative humidity, unless an
  alternative range is explicitly declared.\\
\item
  \textbf{Environmental Sealing:} Conformal coating, potting, or
  physical sealing SHOULD be applied for models explicitly marketed for
  humid, corrosive, or outdoor environments.\\
\item
  Condensation prevention bounds SHALL be documented, and survivability
  under thermal cycling SHOULD be verified.
\end{itemize}

\subsubsection{9.2.3 Altitude and Airflow}\label{altitude-and-airflow}

\begin{itemize}
\tightlist
\item
  Operating altitude range and corresponding thermal derating curves
  SHALL be defined (e.g., typically 0--3000 meters above sea level).\\
\item
  Host airflow assumptions SHALL be documented, including the minimum
  required airflow (e.g., Linear Feet per Minute - LFM or CFM) to
  guarantee safe heat dissipation as a direct function of the inlet air
  temperature.
\end{itemize}

\begin{center}\rule{0.5\linewidth}{0.5pt}\end{center}

\subsection{9.3 Mechanical and Structural
Requirements}\label{mechanical-and-structural-requirements}

\subsubsection{9.3.1 Form Factor}\label{form-factor}

\begin{itemize}
\tightlist
\item
  The physical envelope SHALL be documented and SHOULD conform to an
  industry-standard dimensioning specification (e.g., PCIe low-profile,
  full-height, OCP NIC 3.0, or standard MEZZ).\\
\item
  Custom, proprietary, or deeply embedded modules SHALL explicitly
  document mounting hole patterns, insertion force limits, and
  recommended connector retention torque levels.
\end{itemize}

\subsubsection{9.3.2 Shock and Vibration}\label{shock-and-vibration}

\begin{itemize}
\tightlist
\item
  TimeCards SHOULD withstand mechanical shock and vibration profiles
  aligned with their intended use cases. \emph{(Note: For example,
  standard server-grade profiles may utilize IEC 60068-2-6 and IEC
  60068-2-27 for vibration and shock respectively, while aerospace or
  heavy industrial profiles may dictate far stricter constraints. A
  single shock requirement is not broadly applicable across all TimeCard
  deployment profiles).}\\
\item
  The mechanical isolation of the oscillator from board-level vibration
  is a matter of proprietary design.\\
\item
  The manufacturer SHALL document the maximum frequency stability
  degradation (e.g., acceleration sensitivity, \(\Gamma\), measured in
  parts per \(g\)) under anticipated vibration levels.
\end{itemize}

\subsubsection{9.3.3 Connectors and
Retention}\label{connectors-and-retention}

\begin{itemize}
\tightlist
\item
  RF and timing reference ports SHALL utilize locking, threaded, or
  high-retention physical connectors (e.g., SMA, SMB, MCX, MMCX).\\
\item
  Faceplate physical labeling SHALL clearly identify the function of
  each exposed port (e.g., GNSS, PPS, 10 MHz, ToD, MGMT).\\
\item
  Appropriate physical cable strain relief SHOULD be incorporated to
  prevent mechanical fatigue on critical RF solder joints.
\end{itemize}

\begin{center}\rule{0.5\linewidth}{0.5pt}\end{center}

\subsection{9.4 Electrical and Power
Environment}\label{electrical-and-power-environment}

\subsubsection{9.4.1 Power Supply}\label{power-supply}

\begin{itemize}
\tightlist
\item
  TimeCards SHALL distinctly define input power rail nominal voltages
  and required absolute tolerances (e.g., +12 V ±5\%, +3.3 V ±3\%).\\
\item
  Devices physically powered via standard host buses (e.g., PCIe slots)
  SHALL comply rigidly with the bus specifications for power sequencing,
  inrush current, and maximum continuous current draw limits.\\
\item
  External power connectors SHOULD include localized reverse-polarity
  and surge protection logic.\\
\item
  Localized energy storage (e.g., supercapacitors or lithium backup
  batteries) MAY be utilized to sustain RTC or holdover operation during
  brief power outages.
\end{itemize}

\subsubsection{9.4.2 Electromagnetic Compatibility
(EMC)}\label{electromagnetic-compatibility-emc}

\begin{itemize}
\tightlist
\item
  Implementations SHALL meet the EMC emission and immunity requirements
  dictated by the destination market's regulatory bodies (e.g., EN
  55032, EN 55035, FCC Part 15 Subpart B).\\
\item
  Baseboard shielding, ground-plane flooding, and oscillator enclosure
  isolation to exclude EMI are matters of proprietary vendor design.
\end{itemize}

\subsubsection{9.4.3 Electrostatic Discharge
(ESD)}\label{electrostatic-discharge-esd}

\begin{itemize}
\tightlist
\item
  ESD protection mechanisms SHALL be provided on all externally exposed
  connectors in accordance with relevant electrical standards (e.g., IEC
  61000-4-2 limits for contact and air discharge).\\
\item
  Safe handling procedures and static warning labels SHOULD be included
  in both manufacturing integration documentation and end-user manuals.
\end{itemize}

\begin{center}\rule{0.5\linewidth}{0.5pt}\end{center}

\subsection{9.5 Environmental Qualification
Tests}\label{environmental-qualification-tests}

All TimeCards intended for commercial production SHOULD undergo a
formalized qualification testing matrix. Vendors MAY utilize the
following baseline parameters, or add validated procedures specific to
their target market:

{\def\LTcaptype{none} % do not increment counter
\begin{longtable}[]{@{}
  >{\raggedright\arraybackslash}p{(\linewidth - 4\tabcolsep) * \real{0.2222}}
  >{\raggedright\arraybackslash}p{(\linewidth - 4\tabcolsep) * \real{0.4074}}
  >{\raggedright\arraybackslash}p{(\linewidth - 4\tabcolsep) * \real{0.3704}}@{}}
\toprule\noalign{}
\begin{minipage}[b]{\linewidth}\raggedright
Test Profile
\end{minipage} & \begin{minipage}[b]{\linewidth}\raggedright
Reference Standard
\end{minipage} & \begin{minipage}[b]{\linewidth}\raggedright
Primary Purpose
\end{minipage} \\
\midrule\noalign{}
\endhead
\bottomrule\noalign{}
\endlastfoot
\textbf{Thermal Cycling} & IEC 60068-2-14 & Validate oscillator phase
stability across temperature extremes. \\
\textbf{Humidity Endurance} & IEC 60068-2-78 & Assess component
corrosion limits and moisture protection. \\
\textbf{Mechanical Shock} & IEC 60068-2-27 & Verify PCB fracture
survivability and mechanical retention. \\
\textbf{Random Vibration} & IEC 60068-2-64 & Characterize phase noise
and frequency stability under active vibration. \\
\textbf{ESD Immunity} & IEC 61000-4-2 & Verify localized protection
against electrostatic discharge. \\
\textbf{EMC Emission/Immunity} & EN 55032 / EN 55035 & Baseline
regulatory compliance. \\
\textbf{Power Interruption} & IEC 61000-4-11 & Evaluate state-machine
restart and holdover recovery transitions. \\
\end{longtable}
}

Performance deviation limits recorded during post-qualification testing
SHALL NOT exceed the vendor's formally declared tolerances.

\begin{center}\rule{0.5\linewidth}{0.5pt}\end{center}

\subsection{9.6 Reliability and
Lifetime}\label{reliability-and-lifetime}

\subsubsection{9.6.1 Mean Time Between Failures
(MTBF)}\label{mean-time-between-failures-mtbf}

\begin{itemize}
\tightlist
\item
  Vendors SHALL specify an expected MTBF utilizing recognized
  reliability modeling standard (e.g., Telcordia SR-332, MIL-HDBK-217F,
  or field-measured historical data).\\
\item
  MTBF goals SHOULD be scaled to the target environment (e.g.,
  \textgreater100,000 hours for standard datacenters, or
  \textgreater50,000 hours for harsh field conditions).
\end{itemize}

\subsubsection{9.6.2 Wear-Out and Service
Life}\label{wear-out-and-service-life}

\begin{itemize}
\tightlist
\item
  Hardware components susceptible to finite lifespans or wear-out (e.g.,
  lithium backup batteries, localized cooling fans, flash memory write
  cycles) SHALL distinctly list their rated service intervals in the
  product documentation.\\
\item
  Essential internal device configuration and persistent calibration
  data SHALL safely persist across expected power cycles over the
  device's operational lifetime.
\end{itemize}

\subsubsection{9.6.3 Calibration and
Traceability}\label{calibration-and-traceability}

\begin{itemize}
\tightlist
\item
  TimeCards SHOULD undergo individual factory calibration against a
  traceable primary standard (e.g., UTC(NPL), UTC(NIST)) before
  delivery.\\
\item
  Available calibration certificates SHOULD record the date of
  calibration, the strict environmental conditions present during the
  test, and the calculated margins of uncertainty.
\end{itemize}

\begin{center}\rule{0.5\linewidth}{0.5pt}\end{center}

\section{Annex A (Informative): Supplemental Compliance and
Documentation
Guidelines}\label{annex-a-informative-supplemental-compliance-and-documentation-guidelines}

\emph{Editor's Note: The contents of this Annex are strictly
informative. Hardware vendors are legally bound to comply with the
regulatory constraints, safety laws, and environmental policies of the
jurisdictions in which their products are manufactured and sold. This
standard aims to aid interoperability, but conformity to this standard
does not implicitly grant or enforce legal regulatory compliance.}

\subsection{A.1 Safety Standards}\label{a.1-safety-standards}

\begin{itemize}
\tightlist
\item
  TimeCard hardware is commonly evaluated against general electrical
  safety standards such as IEC 62368-1.
\item
  Good engineering practices encourage the incorporation of
  over-voltage, over-current, and thermal shutdown thermal runaway
  protections.
\item
  Required safety labels, high-voltage warnings, and chassis grounding
  instructions are typically made visible on the physical unit per
  market law.
\end{itemize}

\subsection{A.2 Environmental
Compliance}\label{a.2-environmental-compliance}

\begin{itemize}
\tightlist
\item
  Products sold into global markets generally conform to hazardous
  material directives such as RoHS, REACH, and WEEE.
\item
  Hardware manufacturers are widely encouraged to utilize halogen-free
  PCB materials and design for extensive end-of-life recyclability.
\end{itemize}

\subsection{A.3 End-of-Life Handling and
Disposal}\label{a.3-end-of-life-handling-and-disposal}

\begin{itemize}
\tightlist
\item
  It is recommended that vendors provide guidance for environmentally
  responsible hardware disposal.
\item
  Best practices dictate that hazardous materials (e.g., lithium coin
  cells or backup batteries) are physically removed and sorted prior to
  e-waste recycling.
\item
  A secure mechanism for the cryptographic erasure of firmware, private
  keys, and proprietary calibration data can significantly aid the
  secure decommissioning of units deployed in highly sensitive
  infrastructure.
\end{itemize}

\subsection{A.4 Suggested Datasheet
Inclusions}\label{a.4-suggested-datasheet-inclusions}

Comprehensive technical datasheets and integration manuals significantly
reduce friction for end-users. Manufacturers are encouraged to document
the following openly: - Complete environmental operating limits and a
summary of qualification testing results. - Calculated MTBF and
reliability datasets. - Proof of ESD and EMC compliance certifications.
- Required power, thermal dissipation (BTU/hr), and CFM airflow
requirements. - Publicly available API/Register mappings to support
unhindered multi-vendor interoperability.

\section{10. Applications and Best Practices
(Informative)}\label{applications-and-best-practices-informative}

This chapter provides informative guidance regarding key deployment
scenarios, application domains, and operational best practices for
integrating and maintaining TimeCard systems. It is intended to assist
system architects, integrators, and facility operators in achieving
optimal synchronization accuracy, stability, and reliability across
diverse environments.

As an informative chapter, the guidelines presented herein are
recommendations and do not constitute normative requirements for IEEE
P3335 conformance.

\begin{center}\rule{0.5\linewidth}{0.5pt}\end{center}

\subsection{10.1 Overview}\label{overview-4}

TimeCards function as precision timing subsystems capable of providing
highly stable, traceable, and interoperable time and frequency services
to host platforms. Because the core architecture abstracts the
complexities of hardware-level synchronization away from the host CPU,
TimeCards are versatile enough to be deployed across a vast array of
industries.

Adopting rigorous best practices for physical installation, software
configuration, and continuous monitoring greatly increases the
likelihood that a TimeCard will operate reliably at the peak of its
specified performance envelope throughout its operational lifecycle.

\begin{center}\rule{0.5\linewidth}{0.5pt}\end{center}

\subsection{10.2 Application Domains}\label{application-domains}

The following subclauses detail the primary industries that leverage
TimeCard architectures, noting the specific operational priorities and
recommended configurations for each.

\subsubsection{10.2.1 Data Centers and Cloud
Infrastructure}\label{data-centers-and-cloud-infrastructure}

\begin{itemize}
\tightlist
\item
  \textbf{Operational Purpose:} Coordinating distributed databases,
  globally synchronized transactions (e.g., Google Spanner/TrueTime
  equivalents), AI training cluster synchronization, and event logging
  across massive computing fleets.
\item
  \textbf{Key Architectures:} PTP networks acting as the primary
  distribution method across the datacenter fabric, with TimeCards
  acting as Grandmaster clocks or high-precision Ordinary Clocks at the
  top-of-rack (ToR).
\item
  \textbf{Recommended Best Practice:} Implement redundant TimeCards
  across multiple racks. Utilize PCIe Precision Time Measurement (PTM)
  to deterministically bridge the time from the network interface card
  directly to the host CPU domain, targeting sub-microsecond absolute
  accuracy.
\end{itemize}

\subsubsection{10.2.2 Telecommunications and 5G/6G
Networks}\label{telecommunications-and-5g6g-networks}

\begin{itemize}
\tightlist
\item
  \textbf{Operational Purpose:} Providing strict phase and frequency
  alignment for Open RAN (O-RAN) baseband units, fronthaul/backhaul
  cellular networks, and edge computing synchronization to prevent
  cellular interference.
\item
  \textbf{Key Architectures:} Utilizing specialized telecom profiles
  such as ITU-T G.8275.1 (full timing support from the network) and
  G.8275.2 (partial timing support).
\item
  \textbf{Recommended Best Practice:} Deploy a comprehensive ensemble of
  reference sources (e.g., GNSS augmented by network-delivered PTP).
  Operators are encouraged to actively monitor holdover Maximum Time
  Interval Error (MTIE) telemetry to anticipate network degradation
  during GNSS spoofing or jamming events.
\end{itemize}

\subsubsection{10.2.3 Financial Systems and High-Frequency Trading
(HFT)}\label{financial-systems-and-high-frequency-trading-hft}

\begin{itemize}
\tightlist
\item
  \textbf{Operational Purpose:} Enabling ultra-precise hardware
  timestamping for trading events to satisfy stringent regulatory
  compliance frameworks (e.g., MiFID II in Europe, SEC Rule 613 in the
  USA).
\item
  \textbf{Key Architectures:} Direct 1PPS electrical signal distribution
  alongside high-frequency PTP broadcast networks, heavily reliant on
  hardware-timestamping at the exact point of ingress/egress.
\item
  \textbf{Recommended Best Practice:} Maintain strict, mathematically
  provable traceability to Coordinated Universal Time (UTC) utilizing
  GNSS-disciplined master clocks. Calibration certificates SHOULD be
  maintained and refreshed annually to satisfy regulatory audits.
\end{itemize}

\subsubsection{10.2.4 Power Grid and Industrial Control Systems
(ICS)}\label{power-grid-and-industrial-control-systems-ics}

\begin{itemize}
\tightlist
\item
  \textbf{Operational Purpose:} Synchronizing distributed control
  systems, phasor measurement units (PMUs), SCADA networks, and
  high-voltage protection relays to monitor wide-area grid stability.
\item
  \textbf{Key Architectures:} Utilizing the IEEE C37.238 PTP Power
  Profile over ruggedized deterministic Ethernet.
\item
  \textbf{Recommended Best Practice:} Focus on environmental hardening.
  TimeCards deployed in substations typically require extended operating
  temperature ranges (−40°C to +85°C) and robust electromagnetic
  interference (EMI) shielding to survive high-voltage switching
  transients.
\end{itemize}

\subsubsection{10.2.5 Scientific Research and
Metrology}\label{scientific-research-and-metrology}

\begin{itemize}
\tightlist
\item
  \textbf{Operational Purpose:} Distributing phase-coherent time
  signatures across particle accelerators, radio telescope arrays, and
  advanced metrology laboratories.
\item
  \textbf{Key Architectures:} Leveraging White Rabbit (WR) optical
  networks or specialized continuous-wave RF distribution (e.g., 10 MHz
  sine waves) to achieve sub-nanosecond precision.
\item
  \textbf{Recommended Best Practice:} Utilize high-stability local
  oscillators (such as Rubidium atomic clocks or high-end OCXOs).
  Operators are encouraged to continuously archive Allan Deviation
  (ADEV) and Time Deviation (TDEV) metrics to establish long-term
  baselines for experiment validation.
\end{itemize}

\begin{center}\rule{0.5\linewidth}{0.5pt}\end{center}

\subsection{10.3 Deployment Best
Practices}\label{deployment-best-practices}

\subsubsection{10.3.1 Hardware and Physical
Installation}\label{hardware-and-physical-installation}

\begin{itemize}
\tightlist
\item
  \textbf{Thermal Management:} Precision oscillators are highly
  sensitive to thermal gradients. Installers are advised to verify that
  the host platform provides adequate, steady airflow across the
  TimeCard and to avoid placing the card adjacent to
  high-heat-dissipating components (e.g., flagship GPUs) un-baffled.
\item
  \textbf{Cabling Quality:} For RF inputs and PPS outputs, operators are
  encouraged to utilize phase-stable, heavily shielded coaxial cables.
  Connectors (e.g., SMA) are best torqued to the manufacturer's exact
  specifications to maintain a consistent 50 Ω impedance boundary.
\item
  \textbf{GNSS Antennas:} Route antenna cabling well away from
  high-noise switching power supplies or motorized equipment to minimize
  EMI ingress that degrades the signal-to-noise ratio.
\end{itemize}

\subsubsection{10.3.2 Configuration of Synchronization
Hierarchy}\label{configuration-of-synchronization-hierarchy}

\begin{itemize}
\tightlist
\item
  \textbf{Source Prioritization:} Administrators are advised to
  configure a distinct reference hierarchy based on precision,
  stability, and availability (e.g., defining Local GNSS as Primary,
  Network PTP as Secondary, and adjacent 1PPS as Tertiary).
\item
  \textbf{Ensemble Algorithms:} If the TimeCard supports multi-input
  ensemble modes, enabling this feature allows the card to
  mathematically weigh multiple high-quality references simultaneously,
  smoothing out transient anomalies from any single source.
\item
  \textbf{Holdover Tuning:} Carefully tune holdover operational
  thresholds based on the characterized stability of the specific local
  oscillator model deployed, preventing the TimeCard from serving
  degraded time for too long after a reference loss.
\end{itemize}

\begin{center}\rule{0.5\linewidth}{0.5pt}\end{center}

\subsection{10.4 Operational and Maintenance Best
Practices}\label{operational-and-maintenance-best-practices}

\subsubsection{10.4.1 Monitoring, Telemetry, and
Alarming}\label{monitoring-telemetry-and-alarming}

\begin{itemize}
\tightlist
\item
  \textbf{Continuous Observation:} Leverage out-of-band management
  interfaces (like IPMI, NC-SI, or dedicated I2C/SMBus bridges) to
  continuously poll synchronization health without burdening the host
  CPU.
\item
  \textbf{Log Aggregation:} Stream frequency and phase offset telemetry
  into centralized time-series databases. This allows for the historical
  analysis of oscillator aging and network-induced jitter.
\item
  \textbf{Threshold Alarms:} Configure proactive SNMP traps or Redfish
  events for critical state changes, such as reference loss (LOS), entry
  into holdover states, or sudden internal thermal spikes.
\end{itemize}

\subsubsection{10.4.2 Firmware Lifecycle and
Security}\label{firmware-lifecycle-and-security}

\begin{itemize}
\tightlist
\item
  \textbf{Cryptographic Verification:} Administrators are strongly
  advised to deploy only firmware packages that are cryptographically
  signed by the original hardware vendor.
\item
  \textbf{Secure Staging:} Test all firmware updates on a staging
  TimeCard mirroring the production environment before executing
  fleet-wide deployments.
\item
  \textbf{Access Control:} Restrict remote network access (e.g.,
  REST/gRPC configuration ports) to isolated management VLANs, and
  routinely rotate authentication credentials.
\end{itemize}

\subsubsection{10.4.3 Calibration and Traceability
Maintenance}\label{calibration-and-traceability-maintenance}

\begin{itemize}
\tightlist
\item
  \textbf{Routine Recalibration:} High-end OCXOs and atomic sources
  experience natural, slow frequency drift over years of operation
  (aging). Recalibrating the hardware against a primary standard (e.g.,
  UTC(NIST)) minimizes this accumulated error.
\item
  \textbf{Traceability Documentation:} Preserve and catalog the factory
  calibration certificates, as well as the logs of operational reference
  links, to simplify compliance audits in regulated industries.
\end{itemize}

\begin{center}\rule{0.5\linewidth}{0.5pt}\end{center}

\subsection{10.5 Common Pitfalls and Mitigation
Strategies}\label{common-pitfalls-and-mitigation-strategies}

The following table outlines frequent deployment issues and recommended
architectural or operational responses:

{\def\LTcaptype{none} % do not increment counter
\begin{longtable}[]{@{}
  >{\raggedright\arraybackslash}p{(\linewidth - 4\tabcolsep) * \real{0.1778}}
  >{\raggedright\arraybackslash}p{(\linewidth - 4\tabcolsep) * \real{0.2889}}
  >{\raggedright\arraybackslash}p{(\linewidth - 4\tabcolsep) * \real{0.5333}}@{}}
\toprule\noalign{}
\begin{minipage}[b]{\linewidth}\raggedright
Issue / Symptom
\end{minipage} & \begin{minipage}[b]{\linewidth}\raggedright
Likely Root Cause
\end{minipage} & \begin{minipage}[b]{\linewidth}\raggedright
Recommended Mitigation Strategy
\end{minipage} \\
\midrule\noalign{}
\endhead
\bottomrule\noalign{}
\endlastfoot
\textbf{Constant 1PPS Misalignment} & Reverse polarity settings, or
significant cable-length propagation delay. & Verify rising-edge
vs.~falling-edge configurations. Calibrate for cable propagation delay
(\textasciitilde5 ns per meter). \\
\textbf{Intermittent Reference Loss} & GNSS antenna sky-view
obstruction, or localized RF jamming/spoofing. & Implement a secondary
network-delivered PTP reference as a fallback. Employ anti-spoofing GNSS
receivers. \\
\textbf{Excessive Drift in Holdover} & Drastic thermal fluctuations
affecting the oscillator during the holdover period. & Improve host
chassis thermal regulation. Utilize a TimeCard model equipped with a
higher-grade OCXO. \\
\textbf{High Host-to-Card Latency} & Relying on software interrupts and
unoptimized OS networking stacks to fetch time. & Transition to in-band
hardware timestamping mechanisms such as PCIe PTM to bypass OS
scheduling jitter. \\
\textbf{Management Interface Timeout} & Host CPU exhaustion preventing
the OS from responding to in-band telemetry polls. & Shift telemetry
polling to out-of-band (OOB) BMC channels like SMBus or NC-SI that
operate independently of the host OS. \\
\end{longtable}
}

\begin{center}\rule{0.5\linewidth}{0.5pt}\end{center}

\subsection{10.6 Summary}\label{summary-2}

By implementing the engineering, integration, and operational best
practices detailed in this chapter, system architects mitigate the risks
associated with distributing highly precise time across complex
distributed systems.

These informative recommendations complement the normative hardware and
logical requirements defined within the broader IEEE P3335
specification. When applied comprehensively, they facilitate the
deployment of robust, resilient, and highly traceable TimeCard
infrastructures that perform reliably across decades of service.

\begin{center}\rule{0.5\linewidth}{0.5pt}\end{center}

\section{Annex A: Metrics
(Informative)}\label{annex-a-metrics-informative}

This annex provides definitions, mathematical measurement methodologies,
and interpretation guidelines for the key timing and frequency metrics
utilized throughout the IEEE P3335 TimeCard specification. Because this
is an informative annex, the metrics and recommendations presented
herein do not constitute strict normative requirements, but rather
establish a common engineering vocabulary to facilitate consistent
evaluation of stability, accuracy, and performance across different
hardware implementations and test environments.

All metrics referenced in this annex heavily align with established
industry metrology literature, specifically ITU-T Recommendations G.810
and G.8260, as well as IEEE Std 1139.

\begin{center}\rule{0.5\linewidth}{0.5pt}\end{center}

\subsection{A.1 Overview}\label{a.1-overview}

Metrology metrics quantify the operational performance of TimeCards
across four primary domains: 1. \textbf{Time Stability:} The long-term
coherence of the clock's phase. 2. \textbf{Frequency Stability:} The
consistency of the clock's oscillator over specified averaging periods.
3. \textbf{Phase Noise:} The short-term, frequency-domain spectral
purity of the oscillator. 4. \textbf{Holdover Behavior:} The predictable
degradation of the timescale when external references are lost.

Utilizing a standardized mathematical framework allows operators to
confidently compare devices from competing vendors and validates
interoperability across heterogeneous networks.

\begin{center}\rule{0.5\linewidth}{0.5pt}\end{center}

\subsection{A.2 Core Timing and Stability
Metrics}\label{a.2-core-timing-and-stability-metrics}

\subsubsection{A.2.1 Allan Deviation
(ADEV)}\label{a.2.1-allan-deviation-adev}

\textbf{Definition:} A statistical measure characterizing the fractional
frequency stability of an oscillator or a timing system as a function of
the averaging time (\(\tau\)). Unlike standard variance, ADEV converges
for common oscillator noise types (such as flicker frequency noise).\\
\textbf{Formula:}

\[ \sigma_y(\tau) = \sqrt{ \frac{1}{2(N-1)} \sum_{i=1}^{N-1} (\bar{y}_{i+1} - \bar{y}_i)^2 } \]

where \(\bar{y}_i\) represents the continuous, successive fractional
frequency averages measured over intervals of duration \(\tau\), and
\(N\) represents the total number of frequency samples.

\textbf{Purpose:} ADEV is the primary metric for evaluating both
short-term and long-term oscillator frequency stability.\\
\textbf{Reference:} ITU-T G.810 Annex I; IEEE Std 1139.

\textbf{Typical Baselines for Reference:} \textbar{} Oscillator
Technology \textbar{} Typical ADEV (\(\tau = 1\) s) \textbar{} Typical
Application \textbar{}
\textbar------------------\textbar----------------\textbar----------------\textbar{}
\textbar{} TCXO (Temperature-Compensated) \textbar{}
\(1 \times 10^{-9}\) \textbar{} Cost-sensitive edge compute \textbar{}
\textbar{} OCXO (Oven-Controlled) \textbar{} \(1 \times 10^{-11}\)
\textbar{} Datacenter / Telecom boundary clocks \textbar{} \textbar{}
Rubidium (Atomic) \textbar{} \(1 \times 10^{-12}\) \textbar{} Primary
Reference Clocks (PRC) / Core \textbar{} \textbar{} CSAC (Chip-Scale
Atomic) \textbar{} \(5 \times 10^{-12}\) \textbar{} Portable /
SWaP-constrained field units \textbar{}

\subsubsection{A.2.2 Time Deviation
(TDEV)}\label{a.2.2-time-deviation-tdev}

\textbf{Definition:} The time-domain equivalent of ADEV. It represents
the measure of time stability related to the phase variations of the
signal.\\
\textbf{Formula:}

\[ \text{TDEV}(\tau) = \frac{\tau}{\sqrt{3}} \times \text{Modified ADEV}(\tau) \]

\textbf{Purpose:} TDEV specifically quantifies the variation in absolute
time error over a given observation period.\\
\textbf{Reference:} ITU-T G.810 Annex I.\\
\textbf{Usage:} TDEV is highly useful for evaluating the precision of
synchronization networks, particularly assessing packet delay variation
(PDV) noise in PTP deployments.

\subsubsection{A.2.3 Maximum Time Interval Error
(MTIE)}\label{a.2.3-maximum-time-interval-error-mtie}

\textbf{Definition:} The maximum peak-to-peak phase or time deviation
observed between any two measurement points within an observation window
of duration \(\tau\).\\
\textbf{Formula:}

\[ \text{MTIE}(\tau) = \max_{i,j} |x_j - x_i|, \quad \text{for all } (t_j - t_i) \le \tau \]

where \(x\) is the time error constraint.

\textbf{Purpose:} MTIE acts as a worst-case boundary metric, measuring
the absolute maximum wander or drift of the timescale.\\
\textbf{Reference:} ITU-T G.8260 Appendix II.5.\\
\textbf{Usage:} MTIE is universally utilized to define strict holdover
performance boundaries (e.g., assessing if a clock drifts more than
\(1.5\) \(\mu s\) over a 24-hour holdover period).

\begin{center}\rule{0.5\linewidth}{0.5pt}\end{center}

\subsection{A.3 Frequency Domain and Jitter
Metrics}\label{a.3-frequency-domain-and-jitter-metrics}

\subsubsection{\texorpdfstring{A.3.1 Phase Noise
(\(\mathcal{L}(f)\))}{A.3.1 Phase Noise (\textbackslash mathcal\{L\}(f))}}\label{a.3.1-phase-noise-mathcallf}

\textbf{Definition:} The power spectral density of phase fluctuations of
a continuous periodic signal, traditionally expressed in decibels
relative to the carrier per hertz (\(\text{dBc/Hz}\)) at varying offset
frequencies from the main carrier frequency.\\
\textbf{Formula:}

\[ \mathcal{L}(f) = 10 \log_{10}\left( \frac{S_\phi(f)}{2\text{ rad}^2} \right) \]

where \(S_\phi(f)\) is the one-sided spectral density of the phase
deviations.

\textbf{Purpose:} Phase noise evaluates the short-term spectral purity
of the local oscillator.\\
\textbf{Reporting Profile:} Manufacturers typically plot phase noise
across logarithmic decades: \(1\text{ Hz}\), \(10\text{ Hz}\),
\(100\text{ Hz}\), \(1\text{ kHz}\), \(10\text{ kHz}\), and
\(100\text{ kHz}\) offsets.\\
\textbf{Reference:} IEEE Std 1139.

\subsubsection{A.3.2 Time Jitter}\label{a.3.2-time-jitter}

\textbf{Definition:} The short-term variance or instability of a
time-domain signal (such as a 1PPS edge or a 10 MHz square wave) from
its ideal, mathematically perfect periodic position.\\
\textbf{Measurement Types:} - \textbf{Peak-to-Peak Jitter
(\(J_{pk-pk}\)):} The absolute maximum time difference between the
earliest and latest occurrence of the signal edge over the sample set. -
\textbf{RMS Jitter (\(J_{rms}\)):} The root-mean-square average of the
jitter samples, representing the common standard deviation of the timing
error.

\textbf{Measurement Best Practice:} Jitter characterization is typically
captured using calibrated Time Interval Counters (TIC) operating over a
statistically significant dataset (e.g., \(1 \times 10^6\) samples) with
picosecond-level native resolution.

\begin{center}\rule{0.5\linewidth}{0.5pt}\end{center}

\subsection{A.4 Accuracy, Precision, and
Granularity}\label{a.4-accuracy-precision-and-granularity}

To prevent ambiguity in TimeCard documentation, developers are
encouraged to consistently separate the definitions of accuracy,
precision, and hardware granularity.

{\def\LTcaptype{none} % do not increment counter
\begin{longtable}[]{@{}
  >{\raggedright\arraybackslash}p{(\linewidth - 4\tabcolsep) * \real{0.1714}}
  >{\raggedright\arraybackslash}p{(\linewidth - 4\tabcolsep) * \real{0.3714}}
  >{\raggedright\arraybackslash}p{(\linewidth - 4\tabcolsep) * \real{0.4571}}@{}}
\toprule\noalign{}
\begin{minipage}[b]{\linewidth}\raggedright
Metric
\end{minipage} & \begin{minipage}[b]{\linewidth}\raggedright
Contextual Definition
\end{minipage} & \begin{minipage}[b]{\linewidth}\raggedright
Example Target Profile
\end{minipage} \\
\midrule\noalign{}
\endhead
\bottomrule\noalign{}
\endlastfoot
\textbf{Accuracy} & The degree of conformance or offset of the measured
time directly to an absolute primary reference scale (e.g., UTC). &
\(<100\text{ ns}\) deviation from UTC(NIST) target \\
\textbf{Precision} & The internal repeatability or statistical variance
of a measurement operation under completely identical conditions. &
\(<5\text{ ns}\) \(1\sigma\) variance over \(1,\!000\) consecutive PTM
transactions \\
\textbf{Resolution} & The absolute smallest physically distinguishable
change the circuit can detect or measure. & \(10\text{ ps}\) inherent
TDC (Time-to-Digital Converter) physical resolution \\
\textbf{Granularity} & The minimal mathematical discrete step available
in the reported digital timestamps. & \(1\text{ ns}\) bit-level
granularity within an IEEE 1588 packet payload \\
\end{longtable}
}

\begin{center}\rule{0.5\linewidth}{0.5pt}\end{center}

\subsection{A.5 Holdover Metrics}\label{a.5-holdover-metrics}

When a TimeCard suffers a total loss of its inbound external reference
signals (e.g., GNSS jamming or fiber cut), the active disciplining loop
opens, and the unit transitions to \textbf{holdover}.

\subsubsection{A.5.1 Holdover Time
Error}\label{a.5.1-holdover-time-error}

\textbf{Definition:} The continuously accumulating deviation of absolute
time driven by the free-running fractional frequency offset of the
internal oscillator during holdover.\\
\textbf{Mathematical Concept:}

\[ \Delta x(t) = x_0 + \int_0^t \Delta y(\tau) d\tau + \frac{1}{2} D t^2 + \epsilon(t) \]

where \(x_0\) is initial phase error, \(\Delta y\) is the normalized
frequency offset, \(D\) is the linear frequency aging rate, and
\(\epsilon(t)\) represents internal random noise.

\textbf{Practical Example:} If a TimeCard enters holdover with an
uncorrected frequency offset of \(1.15 \times 10^{-11}\), the generated
1PPS signal will mechanically drift by approximately \(1\text{ \mu s}\)
per day of holdover (\(\approx 86.4\text{ ns}\) per
\(1 \times 10^{-12}\) offset).

\subsubsection{A.5.2 Warm-Up and Recovery
Behavior}\label{a.5.2-warm-up-and-recovery-behavior}

\begin{itemize}
\tightlist
\item
  \textbf{Warm-Up Time:} The total duration required for an oscillator
  (specifically heated OCXOs or Rubidium cells) to reach thermal
  equilibrium and electrical stabilization upon a cold power cycle.
\item
  \textbf{Recovery Time:} The architectural time required for the PLL to
  re-establish a solid lock and compress phase errors after successfully
  reacquiring a valid reference.
\item
  \textbf{Note:} It is highly recommended that vendors empirically
  characterize and publish both parameters in standard hardware
  datasheets.
\end{itemize}

\begin{center}\rule{0.5\linewidth}{0.5pt}\end{center}

\subsection{A.6 Ensemble and Correlation
Metrics}\label{a.6-ensemble-and-correlation-metrics}

In distributed environments, multiple TimeCards or multiple reference
inputs may be mathematically combined into an \textbf{Ensemble Clock}.

{\def\LTcaptype{none} % do not increment counter
\begin{longtable}[]{@{}
  >{\raggedright\arraybackslash}p{(\linewidth - 2\tabcolsep) * \real{0.4091}}
  >{\raggedright\arraybackslash}p{(\linewidth - 2\tabcolsep) * \real{0.5909}}@{}}
\toprule\noalign{}
\begin{minipage}[b]{\linewidth}\raggedright
Concept
\end{minipage} & \begin{minipage}[b]{\linewidth}\raggedright
Description
\end{minipage} \\
\midrule\noalign{}
\endhead
\bottomrule\noalign{}
\endlastfoot
\textbf{Ensemble Average Stability} & Statistical improvement in
composite ADEV proportional to \(\sqrt{N}\), where \(N\) represents the
number of completely independent, uncorrelated clocks operating in the
ensemble. \\
\textbf{Correlation Coefficient (\(\rho\))} & A statistical measurement
quantifying the dependency or shared error vectors between clocks. A
highly resilient ensemble relies on sources with low correlation
(\(\rho \approx 0\)). \\
\textbf{Weighting Function} & The active algorithm determining the
proportional contribution of each individual clock/reference step,
dynamically shifting based on real-time MTIE performance or variance. \\
\end{longtable}
}

\begin{center}\rule{0.5\linewidth}{0.5pt}\end{center}

\subsection{A.7 Environmental Influence
Metrics}\label{a.7-environmental-influence-metrics}

Because oscillators physically react to their environment, static
laboratory metrics rarely reflect real-world holdover stability unless
environmental parameters are factored in.

{\def\LTcaptype{none} % do not increment counter
\begin{longtable}[]{@{}
  >{\raggedright\arraybackslash}p{(\linewidth - 4\tabcolsep) * \real{0.2791}}
  >{\raggedright\arraybackslash}p{(\linewidth - 4\tabcolsep) * \real{0.2093}}
  >{\raggedright\arraybackslash}p{(\linewidth - 4\tabcolsep) * \real{0.5116}}@{}}
\toprule\noalign{}
\begin{minipage}[b]{\linewidth}\raggedright
Parameter
\end{minipage} & \begin{minipage}[b]{\linewidth}\raggedright
Metric Definition
\end{minipage} & \begin{minipage}[b]{\linewidth}\raggedright
Impact
\end{minipage} \\
\midrule\noalign{}
\endhead
\bottomrule\noalign{}
\endlastfoot
\textbf{Temperature Coefficient} & The variance in fractional frequency
(\(\Delta f/f\)) per degree Celsius (\(\Delta^\circ\text{C}\)). &
Directly dictates holdover drift if the chassis undergoes drastic
thermal shifts (e.g., an AC failure). \\
\textbf{Vibration Sensitivity (\(\Gamma\))} & The relative frequency
change per \(g\) of mechanical acceleration. & Critical factor for
systems deployed in heavily vibrating racks or industrial transit
enclosures. \\
\textbf{Aging Rate} & The constant, predictable, unidirectional
long-term frequency drift of the resonator over time (days/years). &
Determines how frequently the TimeCard requires absolute external
recalibration to prevent systemic clock bias. \\
\end{longtable}
}

\begin{center}\rule{0.5\linewidth}{0.5pt}\end{center}

\subsection{A.8 Evaluating Hardware: Best
Practices}\label{a.8-evaluating-hardware-best-practices}

When designing lab tests or automated qualification fixtures to evaluate
the metrics detailed in this annex, engineers are encouraged to observe
the following best practices: - \textbf{Traceability:} Utilize only
highly calibrated measurement instruments (e.g., Phase Noise Analyzers,
Counters) that possess verifiable traceability to primary national
standards (NIST, PTB, NPL). - \textbf{Control Symmetries:} Rigorously
design the physical measurement setup to minimize asymmetric cable
propagation delays, thermal gradients, and RF impedance mismatches. -
\textbf{Broad Spectrums:} Execute stability validation tests across
broad averaging intervals (\(\tau = 1\text{ s}\) out to
\(>10^5\text{ s}\)) to capture both short-term jitter and diurnal
temperature-driven wander. - \textbf{Contextual Logging:} Always log
environmental metadata concurrently with performance data (e.g., ambient
temperature, relative humidity, chassis airflow, and test equipment
calibration schedules).

\begin{center}\rule{0.5\linewidth}{0.5pt}\end{center}

\subsection{A.9 Summary}\label{a.9-summary}

The mathematical metrics outlined in this annex establish a vital,
unified vocabulary for characterizing the temporal performance of
TimeCard architectures.

By standardizing exactly how stability, noise, drift, and holdover are
mathematically calculated and reported, the IEEE P3335 ecosystem allows
operators to rapidly compare discrete hardware configurations, validate
regulatory compliance, and confidently integrate mixed-vendor timing
infrastructure.

\begin{center}\rule{0.5\linewidth}{0.5pt}\end{center}

\textbf{End of Annex A}

\section{Annex B: Test Procedures
(Informative)}\label{annex-b-test-procedures-informative}

This annex provides standardized, repeatable engineering test procedures
for evaluating the functional and performance characteristics of
TimeCard implementations. These procedures are designed to ensure
consistency, comparability, and mathematical traceability of measurement
results across independent vendors, metrology laboratories, and live
deployment environments.

As this annex is explicitly informative, the test methodologies proposed
herein serve as best-practice engineering baselines for hardware
validation rather than strict normative conformance requirements.

\begin{center}\rule{0.5\linewidth}{0.5pt}\end{center}

\subsection{B.1 Overview}\label{b.1-overview}

Comprehensive validation of a TimeCard architecture encompasses three
primary testing domains: 1. \textbf{Functional Verification:} Validating
logical compliance with the architectural interface, physical
connectability, and management plane state-machine behaviors. 2.
\textbf{Performance Validation:} Empirically measuring TimeCard
frequency stability, phase noise, holdover drift, and timestamping
jitter against the vendor's declared specifications. 3.
\textbf{Environmental and Reliability Qualification:} Assessing physical
performance consistency and survivability under externally applied,
simulated stress conditions (e.g., thermal shock, vibration, or EMI).

To satisfy traceability requirements, it is strongly recommended that
all tests be conducted utilizing equipment recently calibrated against
recognized national metrology standards.

\begin{center}\rule{0.5\linewidth}{0.5pt}\end{center}

\subsection{B.2 General Test
Conditions}\label{b.2-general-test-conditions}

\subsubsection{B.2.1 Environmental
Setup}\label{b.2.1-environmental-setup}

\begin{itemize}
\tightlist
\item
  Testing environments are typically maintained under nominal, thermally
  stable laboratory conditions unless a specific thermal extreme test is
  actively being executed. Baseline conditions generally sit at
  \(23 \pm 3 ^\circ\text{C}\) with \(30\text{--}70\%\) non-condensing
  relative humidity.
\item
  The Device Under Test (DUT) should be allowed to reach complete
  thermal equilibrium and internal oscillator operational stability
  (warm-up) prior to the commencement of any metrology recording.
\item
  Physical cabling (e.g., coaxial cables for 1PPS or 10 MHz) and RF
  splitters utilized in the test fixture should be phase-matched and
  carefully impedance-matched (typically \(50\,\Omega\)) to prevent
  signal reflections that cause false jitter readings in the
  instrumentation.
\end{itemize}

\subsubsection{B.2.2 Power and
Initialization}\label{b.2.2-power-and-initialization}

\begin{itemize}
\tightlist
\item
  Evaluators should apply host bus power (e.g., via a PCIe riser) or
  auxiliary power within the hardware's rated voltage tolerance.
\item
  The initial power-up sequence provides a baseline to map the
  oscillator's raw warm-up time from cold-start to achieving a
  steady-state frequency lock.
\end{itemize}

\subsubsection{B.2.3 Measurement Equipment Calibration
limits}\label{b.2.3-measurement-equipment-calibration-limits}

Instrumentation utilized in performance validation should possess a
noise floor or error margin substantially lower than the DUT to ensure
the DUT is being measured, rather than the noise of the instrumentation
itself. - \textbf{Reference Clock:} The primary laboratory reference
clock driving the test equipment should possess an Allan Deviation
(ADEV) at least one order of magnitude (\(10\times\)) superior to the
target specification of the DUT. - \textbf{Time Interval Counter (TIC):}
TICs should possess a native single-shot resolution of \(<10\text{ ps}\)
for evaluating high-precision OCXOs or GNSS-disciplined TimeCards. -
\textbf{Phase Noise Analyzer:} The analyzer's internal local oscillator
should possess a phase noise floor at least \(10\text{ dB}\), preferably
\(15\text{ dB}\), lower than the anticipated mask of the DUT across all
relevant frequency offsets.

\begin{center}\rule{0.5\linewidth}{0.5pt}\end{center}

\subsection{B.3 Functional Verification
Tests}\label{b.3-functional-verification-tests}

Functional tests validate the logical operations and state transitions
of the TimeCard without necessarily grading sub-nanosecond physical
precision.

\subsubsection{B.3.1 Receive (Inbound) Interface
Tests}\label{b.3.1-receive-inbound-interface-tests}

\begin{enumerate}
\def\labelenumi{\arabic{enumi}.}
\tightlist
\item
  \textbf{Reference Detection:} Sequentially apply valid external
  references (e.g., GNSS antenna input, PTP network feed, external 1PPS)
  and verify via the management telemetry that the TimeCard correctly
  recognizes the presence of the physical signal.
\item
  \textbf{Lock Acquisition:} Monitor the management plane to record the
  time taken for the TimeCard's internal Phase-Locked Loop (PLL) or
  disciplining mechanism to transition from ``Free-Run'' to ``Locked''
  status. Compare this to the vendor's specified acquisition time limit.
\item
  \textbf{Failover Response:} While actively locked to a primary
  reference, abruptly disconnect the primary input. Verify that the
  hardware deterministically enters the ``Holdover'' state or smoothly
  transitions to a secondary reference without causing a systemic reboot
  or crashing the host driver.
\end{enumerate}

\subsubsection{B.3.2 Providing (Outbound) Interface
Tests}\label{b.3.2-providing-outbound-interface-tests}

\begin{enumerate}
\def\labelenumi{\arabic{enumi}.}
\tightlist
\item
  \textbf{PPS Amplitude and Polarity:} Utilizing an oscilloscope bridged
  with a \(50\,\Omega\) terminator, capture the 1PPS output edge. Verify
  the correct edge polarity (rising vs.~falling), the slew rate (rise
  time), and the high/low voltage thresholds against standard CMOS/LVTTL
  specifications.
\item
  \textbf{PCIe PTM Validation:} If supported, execute host-side software
  to trigger Precision Time Measurement (PTM) dialogs across the PCIe
  bus. Correlate the returned PTM hardware timestamps against the
  physical 1PPS output to verify systemic, in-band latency and offset
  behavior.
\end{enumerate}

\begin{center}\rule{0.5\linewidth}{0.5pt}\end{center}

\subsection{B.4 Performance Validation
Tests}\label{b.4-performance-validation-tests}

Performance tests quantify the physical metrology capabilities of the
TimeCard.

\subsubsection{B.4.1 Frequency Stability
(ADEV/TDEV)}\label{b.4.1-frequency-stability-adevtdev}

\begin{itemize}
\tightlist
\item
  \textbf{Procedure:} Connect the TimeCard's primary continuous clock
  output (e.g., 10 MHz) to a Phase Noise Analyzer or high-resolution
  Frequency Counter. Record continuous phase/frequency samples locked
  against the laboratory primary reference.
\item
  \textbf{Analysis:} Calculate and plot the Allan Deviation
  (\(\sigma_y(\tau)\)) and Time Deviation (TDEV) across logarithmic
  averaging intervals
  (\(\tau = 1\text{ s}, 10\text{ s}, 100\text{ s}, 1,\!000\text{ s}, 10,\!000\text{ s}\)).
\item
  \textbf{Evaluation:} Compare the resulting stability curves against
  the manufacturer's published datasheet limits.
\end{itemize}

\subsubsection{\texorpdfstring{B.4.2 Short-Term Phase Noise
(\(\mathcal{L}(f)\))}{B.4.2 Short-Term Phase Noise (\textbackslash mathcal\{L\}(f))}}\label{b.4.2-short-term-phase-noise-mathcallf}

\begin{itemize}
\tightlist
\item
  \textbf{Procedure:} Utilize a cross-correlated Phase Noise Analyzer to
  plot the Single Sideband (SSB) phase noise of the oscillator.
\item
  \textbf{Analysis:} Measure the power spectral density in
  \(\text{dBc/Hz}\) at standard decade offsets extending from
  \(1\text{ Hz}\) to \(1\text{ MHz}\) from the primary carrier
  frequency.
\item
  \textbf{Evaluation:} Validate that the plotted noise curve sits below
  the vendor's declared upper-bound phase noise mask.
\end{itemize}

\subsubsection{B.4.3 Time Jitter}\label{b.4.3-time-jitter}

\begin{itemize}
\tightlist
\item
  \textbf{Procedure:} Capture a statistically significant sequence of
  1PPS pulses (e.g., minimum \(1 \times 10^5\) samples) utilizing a TIC
  triggered against the pristine laboratory 1PPS reference.
\item
  \textbf{Analysis:} Calculate the Root-Mean-Square (RMS) jitter and the
  absolute peak-to-peak (\(J_{pk-pk}\)) timing variance of the signal
  edges.
\end{itemize}

\subsubsection{B.4.4 Holdover Boundary
(MTIE)}\label{b.4.4-holdover-boundary-mtie}

\begin{itemize}
\tightlist
\item
  \textbf{Procedure:} Allow the DUT to lock to a master reference and
  achieve thermal and mathematical equilibrium for at least 24 hours.
  Abruptly sever the reference connection, forcing the unit into
  holdover.
\item
  \textbf{Analysis:} Continuously record the wandering time error of the
  TimeCard's 1PPS output against the laboratory master 1PPS for the
  specified holdover duration (e.g., 4, 12, or 24 hours).
\item
  \textbf{Evaluation:} Compute the Maximum Time Interval Error (MTIE)
  across the observation window and verify the drift profile remains
  beneath the target limit (e.g., \(1.5\text{ \mu s}\)/24hr).
\end{itemize}

\begin{center}\rule{0.5\linewidth}{0.5pt}\end{center}

\subsection{B.5 Environmental and Stress
Qualification}\label{b.5-environmental-and-stress-qualification}

Environmental tests validate that the TimeCard's hardware design
robustly compensates for external physical interference.

{\def\LTcaptype{none} % do not increment counter
\begin{longtable}[]{@{}
  >{\raggedright\arraybackslash}p{(\linewidth - 4\tabcolsep) * \real{0.1714}}
  >{\raggedright\arraybackslash}p{(\linewidth - 4\tabcolsep) * \real{0.4000}}
  >{\raggedright\arraybackslash}p{(\linewidth - 4\tabcolsep) * \real{0.4286}}@{}}
\toprule\noalign{}
\begin{minipage}[b]{\linewidth}\raggedright
Test Profile
\end{minipage} & \begin{minipage}[b]{\linewidth}\raggedright
Recommended Methodology
\end{minipage} & \begin{minipage}[b]{\linewidth}\raggedright
Verification Target
\end{minipage} \\
\midrule\noalign{}
\endhead
\bottomrule\noalign{}
\endlastfoot
\textbf{Thermal Cycling} & Place the active, locked DUT inside an
environmental chamber. Cycle the ambient temperature across the
manufacturer's maximum specified operating range (e.g.,
\(0^\circ\text{C}\) to \(55^\circ\text{C}\)) using a predefined ramp
rate (\(^\circ\text{C}\)/min). & Assess the maximum frequency excursion
during temperature gradients. Verify that the oscillator's internal
temperature compensation mechanism functions correctly. \\
\textbf{Power Interruption} & Mechanically or electrically interrupt the
host power plane for brief intervals (\(1\text{ s}\) to \(10\text{ s}\))
if the TimeCard relies on localized holdover capacitors or batteries. &
Verify that the hardware successfully bridges the interruption and
maintains the internal unified timescale without experiencing a phase
step or discontinuity upon power restoration. \\
\textbf{Vibration / Shock} & Subject the TimeCard to variable frequency,
multi-axis harmonic vibration (\(5\text{ Hz}\) to \(500\text{ Hz}\))
utilizing a specialized shaker table. & Record the immediate degradation
in phase noise and frequency stability caused by g-force acceleration
(\(\Gamma\)), verifying that the structural design adequately dampens
the oscillator. \\
\end{longtable}
}

\begin{center}\rule{0.5\linewidth}{0.5pt}\end{center}

\subsection{B.6 Reporting and
Traceability}\label{b.6-reporting-and-traceability}

To maintain engineering transparency, comprehensive test reports for
TimeCard validations should document: - A detailed block diagram
describing the physical test fixture, including cable lengths to account
for fixed propagation delays. - Expiration dates and calibration
certificates for all primary laboratory measurement instrumentation used
during the test. - The precise firmware revision, hardware board step,
and baseboard management controller (BMC) version installed on the DUT.
- Raw and plotted measurement data for all ADEV, TDEV, Phase Noise, and
MTIE evaluations. - Explicitly stated conditions during holdover tests
(e.g., the thermal ramp rate applied to the chassis while the reference
was disconnected).

By adhering to a consistent, traceable testing methodology, integrators
can confidently benchmark TimeCard implementations across diverse vendor
ecosystems, guaranteeing reliable behavior in production infrastructure.

\begin{center}\rule{0.5\linewidth}{0.5pt}\end{center}

\textbf{End of Annex B}

\section{Annex C: Bibliography
(Informative)}\label{annex-c-bibliography-informative}

This bibliography lists the key standards, engineering references,
academic publications, and technical documents that provide foundational
knowledge regarding the design, testing, and implementation of TimeCard
architectures.

These documents are strictly non-normative. They are provided for
informational purposes only, and their implementation is not strictly
required to claim structural conformance to the IEEE P3335 standard.

\begin{center}\rule{0.5\linewidth}{0.5pt}\end{center}

\subsection{C.1 International Standards and Industry
Profiles}\label{c.1-international-standards-and-industry-profiles}

\begin{itemize}
\tightlist
\item
  \textbf{{[}B1{]}} CENELEC EN 55032 / EN 55035, \emph{Electromagnetic
  compatibility of multimedia equipment -- Emission and Immunity
  Requirements}, 2017.
\item
  \textbf{{[}B2{]}} DMTF DSP0236, \emph{Intelligent Platform Management
  Interface (IPMI) Specification v2.0}, 2015.
\item
  \textbf{{[}B3{]}} IEC 60068-2-6, \emph{Environmental testing -- Part
  2-6: Tests - Test Fc: Vibration (sinusoidal)}, 2007.
\item
  \textbf{{[}B4{]}} IEC 60068-2-27, \emph{Environmental testing -- Part
  2-27: Tests - Test Ea and guidance: Shock}, 2008.
\item
  \textbf{{[}B5{]}} IEC 61000-4-2, \emph{Electromagnetic compatibility
  (EMC) -- Part 4-2: Testing and measurement techniques - Electrostatic
  discharge immunity test}, 2009.
\item
  \textbf{{[}B6{]}} IEC 62368-1, \emph{Audio/video, information and
  communication technology equipment -- Part 1: Safety requirements},
  2018.
\item
  \textbf{{[}B7{]}} IEEE Std C37.238™-2017, \emph{IEEE Standard Profile
  for Use of IEEE 1588 Precision Time Protocol in Power System
  Applications}, 2017.
\item
  \textbf{{[}B8{]}} IETF RFC 3411, \emph{An Architecture for Describing
  Simple Network Management Protocol (SNMP) Management Frameworks},
  2002.
\item
  \textbf{{[}B9{]}} ISO/IEC 17025, \emph{General requirements for the
  competence of testing and calibration laboratories}, 2017.
\item
  \textbf{{[}B10{]}} ITU-T Recommendation G.8275.1, \emph{Precision time
  protocol telecom profile for phase/time synchronization with full
  timing support from the network}, 2016.
\item
  \textbf{{[}B11{]}} ITU-T Recommendation G.8275.2, \emph{Precision time
  protocol telecom profile for phase/time synchronization with partial
  timing support from the network}, 2017.
\item
  \textbf{{[}B12{]}} PCI-SIG, \emph{PCI Express Base Specification
  Revision 5.0 (incorporating Precision Time Measurement ECN)}, 2019.
\item
  \textbf{{[}B13{]}} SMBus Forum, \emph{System Management Bus (SMBus)
  Specification Version 3.2}, 2018.
\end{itemize}

\begin{center}\rule{0.5\linewidth}{0.5pt}\end{center}

\subsection{C.2 Academic and Technical Metrology
References}\label{c.2-academic-and-technical-metrology-references}

\begin{itemize}
\tightlist
\item
  \textbf{{[}B14{]}} Allan, D. W., \emph{Time and Frequency
  (Time-Domain) Characterization, Estimation, and Prediction of
  Precision Clocks and Oscillators}, IEEE Transactions on Ultrasonics,
  Ferroelectrics, and Frequency Control, 1987.
\item
  \textbf{{[}B15{]}} Bregni, S., \emph{Measurement of Maximum Time
  Interval Error for Telecommunications Clock Stability
  Characterization}, IEEE Transactions on Instrumentation and
  Measurement, Vol. 45, No.~5, 1996.
\item
  \textbf{{[}B16{]}} Calhoun, J., \emph{Clock Stability Characterization
  for Network Synchronization}, IEEE Communications Surveys \&
  Tutorials, 2009.
\item
  \textbf{{[}B17{]}} CERN / GSI, \emph{White Rabbit Specification and
  Documentation (Sub-nanosecond Ethernet synchronization)}.
\item
  \textbf{{[}B18{]}} Department of Defense, \emph{MIL-HDBK-217F:
  Reliability Prediction of Electronic Equipment}, 1991.
\item
  \textbf{{[}B19{]}} Hartmann, J., \emph{Open Timing Architectures for
  Data Centers}, OCP Time Appliances Project Whitepaper, 2020.
\item
  \textbf{{[}B20{]}} Lombardi, M. A., \emph{NIST Technical Note 1952:
  Introduction to PTP and synchronization techniques}, National
  Institute of Standards and Technology, 2016.
\item
  \textbf{{[}B21{]}} Riley, W. J., \emph{NIST Special Publication 1065:
  Handbook of Frequency Stability Analysis}, National Institute of
  Standards and Technology, 2008.
\item
  \textbf{{[}B22{]}} Weiss, M., \emph{Time Transfer and Synchronization
  for Distributed Systems}, NIST Technical Report, 2013.
\end{itemize}

\begin{center}\rule{0.5\linewidth}{0.5pt}\end{center}

\subsection{C.3 Supporting Open Source and Community
Resources}\label{c.3-supporting-open-source-and-community-resources}

\begin{itemize}
\tightlist
\item
  \textbf{{[}B23{]}} Network Time Foundation, \emph{NTP.org Reference
  Implementation (ntpd) source code and documentation}.
\item
  \textbf{{[}B24{]}} Open Compute Project (OCP) Time Appliances Project
  (TAP), \emph{TimeCard Hardware Schematics and Firmware Repository}.
  {[}\url{https://github.com/opencomputeproject/Time-Appliance-Project}{]}
\item
  \textbf{{[}B25{]}} Open Compute Project (OCP) Time Appliances Project
  (TAP), \emph{Interoperability Whitepaper and Multi-Vendor Testing
  Framework}.
\item
  \textbf{{[}B26{]}} Turba Technologies, \emph{Turba Analytics
  Documentation - AI workload and distributed timing analytics
  platform}.
\end{itemize}

\begin{center}\rule{0.5\linewidth}{0.5pt}\end{center}

\subsection{C.4 Related Military/Defense Specifications
(DID)}\label{c.4-related-militarydefense-specifications-did}

\begin{itemize}
\tightlist
\item
  \textbf{{[}B27{]}} DI-IPSC-81431A, \emph{System/Subsystem
  Specification (SSS) Data Item Description}, 2000.
\item
  \textbf{{[}B28{]}} DI-IPSC-81432A, \emph{System/Subsystem Design
  Description (SSDD) Data Item Description}, 1999.
\item
  \textbf{{[}B29{]}} DI-IPSC-81434A, \emph{Interface Requirements
  Specification (IRS) Data Item Description}, 1999.
\item
  \textbf{{[}B30{]}} DI-IPSC-81436A, \emph{Interface Design Description
  (IDD) Data Item Description}, 1999.
\end{itemize}

\begin{center}\rule{0.5\linewidth}{0.5pt}\end{center}

\textbf{End of Annex C}
